\documentclass[11pt]{article}

\usepackage[a4paper,margin=2.5cm]{geometry}
\usepackage[T1]{fontenc}
\usepackage{lmodern}
\usepackage{amsmath,amssymb,amsthm,mathtools}
\usepackage{microtype}
\usepackage{booktabs}
\usepackage{placeins}
\usepackage{enumitem}
\usepackage{pgfplots}
\pgfplotsset{compat=1.18}
\usepgfplotslibrary{groupplots}
\usepackage[hidelinks]{hyperref}
\usepackage[capitalise,nameinlink]{cleveref}
\usepackage{thmtools}

\newcommand{\CfQuadCert}{0.1380826}
\newcommand{\CfCubicCert}{0.135607}
\newcommand{\CfLinCert}{0.142613}
\newcommand{\CfTCert}{0.1284786}
\newcommand{\FloatCheck}{2.3\cdot 10^{-16}}
\newcommand{\StabCubic}{0.91}
\newcommand{\StabOther}{0.04}
\newcommand{\NDiagMin}{0.140}
\newcommand{\NDiagMax}{0.145}
\newcommand{\NOneMin}{0.166}
\newcommand{\NOneMax}{0.184}
\newcommand{\BetaMin}{1.538}
\newcommand{\BetaMax}{1.552}
\newcommand{\StableCutoff}{0.01}
\newcommand{\RelQuad}{-2.61}
\newcommand{\RelCubic}{+0.53}
\newcommand{\RelLin}{-0.65}
\newcommand{\SensQuadMin}{0.134}
\newcommand{\SensQuadMax}{0.142}
\newcommand{\XLargeInvCubeSmall}{1.6}
\newcommand{\XLargeInvCubeOne}{1.6\cdot 10^{-3}}
\newcommand{\XTwoNinths}{0.0135}
\newcommand{\DeltaRowOneMax}{2.4\cdot 10^{-4}}
\newcommand{\DeltaRowMax}{0.103}
\newcommand{\DeltaRowMaxR}{0.1}
\newcommand{\CfQuad}{0.13808263}
\newcommand{\CfCubic}{0.13560752}
\newcommand{\CfLin}{0.14261331}
\newcommand{\CfT}{0.12847867}
\newcommand{\VaughanCtwo}{0.12847867}
\newcommand{\IConst}{0.831544899909}
\newcommand{\ShapeMin}{0.9938}
\newcommand{\ShapeMax}{1.0011}
\newcommand{\ShapeRows}{66}
\newcommand{\ShapeDevLargest}{1.8\cdot 10^{-4}}
\newcommand{\SelfTest}{4.8\cdot 10^{-10}}
\newcommand{\RerunMaxRelV}{0}
\newcommand{\RhoBound}{1000}


\numberwithin{equation}{section}
\declaretheorem[style=plain,numberwithin=section]{theorem}
\declaretheorem[style=plain,sibling=theorem]{proposition}
\declaretheorem[style=plain,sibling=theorem]{lemma}
\declaretheorem[style=remark,sibling=theorem]{remark}
% Section-qualified hyperref anchors, so that equal numbers in different
% sections (and appendices) get distinct link targets.
\renewcommand{\theHequation}{\theHsection.\arabic{equation}}
\renewcommand{\theHtheorem}{\theHsection.\arabic{theorem}}
\providecommand{\theHproposition}{}\renewcommand{\theHproposition}{\theHsection.\arabic{theorem}}
\providecommand{\theHlemma}{}\renewcommand{\theHlemma}{\theHsection.\arabic{theorem}}
\providecommand{\theHremark}{}\renewcommand{\theHremark}{\theHsection.\arabic{theorem}}
\makeatletter
\providecommand*{\toclevel@theorem}{0}
\providecommand*{\toclevel@proposition}{0}
\providecommand*{\toclevel@lemma}{0}
\providecommand*{\toclevel@remark}{0}
\makeatother

% Notation
\newcommand{\N}{\mathbb{N}}
\newcommand{\Z}{\mathbb{Z}}
\newcommand{\R}{\mathbb{R}}
\newcommand{\one}{\mathbf{1}}
\newcommand{\MT}{\mathrm{MT}}
\newcommand{\rad}{\operatorname{rad}}
\newcommand{\lcm}{\operatorname{lcm}}
\newcommand{\fr}[1]{\{#1\}}
\newcommand{\eps}{\varepsilon}
\newcommand{\hnat}{h^{\natural}}

\title{The sparse variance of squarefree numbers in arithmetic progressions to polynomial moduli: an elementary asymptotic formula}
\author{Write-up of an archived machine-generated proof run\\(provenance in Appendix~\ref{app:provenance})}
\date{}

\begin{document}
\maketitle

\begin{abstract}
For a polynomial $f\in\Z[t]$ of degree $d\ge1$ with $f(t)\ge2$ and $f'(t)\ge1$ for $t\ge1$, consider
\[
V_f(x,y)=\sum_{1\le k\le y}f'(k)\sum_{l=1}^{f(k)}\bigl(Q(x;f(k),l)-g(f(k),l)\,x\bigr)^2,
\]
where $Q(x;q,l)$ counts squarefree $n\le x$ with $n\equiv l\pmod q$ and $g(q,l)$ is the corresponding density.
We show by elementary means that
for each fixed $\delta\in(0,1/9)$, $V_f(x,y)=C_f\,x^{1/2}f(y)^{3/2}\bigl(1+O(x^{-\eta})\bigr)$ uniformly for $x^{7/9+\delta}\le f(y)\le x^{1-\delta}$, with an explicit Euler product $C_f>0$ and $\eta=\min(\delta/12,1/(6d))$; the implied constant depends only on $f$ and $\delta$.
We also give explicit error terms valid throughout $f(y)\le x$, and a version with an exact arithmetic main term.
The proof rests on an exact identity for the variance with respect to a single modulus, a uniform count of pairs of squarefree numbers with a given difference, an Euler--Maclaurin analysis of sawtooth sums $\sum_s B(Z/s^2)$, and an average of a multiplicative function over polynomial values.
We compare the result with exact numerical computations of $V_f$ for $f=k^2+1$, $k^3+2$ and $2k+1$ up to $x=2.56\cdot10^8$.
\end{abstract}

\tableofcontents

\section{Introduction}\label{sec:intro}

\subsection{The problem}

Let $\mu^2(n)$ denote the indicator function of the squarefree positive integers. For real $x\ge1$, $q\in\N$ and $l\in\Z$ put
\[
Q(x;q,l)=\#\{1\le n\le x:\ n\equiv l\ (\mathrm{mod}\ q),\ \mu^2(n)=1\},\qquad Q(x)=Q(x;1,0).
\]
The limit $g(q,l)=\lim_{x\to\infty}Q(x;q,l)/x$ exists (\cref{prop:densities}) and is the density of squarefree numbers in the class $l$ modulo $q$. Its error term and the variance with respect to one modulus are
\[
E(x;q,l)=Q(x;q,l)-g(q,l)\,x,\qquad V(x;q)=\sum_{l=1}^{q}E(x;q,l)^2 .
\]
The Montgomery--Hooley-type variance $V(x,Y)=\sum_{q\le Y}V(x;q)$ over all moduli $q\le Y$ was evaluated asymptotically by Vaughan~\cite{Vaughan2005}; we recall his result below.

This paper treats a \emph{sparse} version, in which the moduli are restricted to the values of a polynomial. Throughout, $f\in\Z[t]$ is a polynomial of degree $d\ge1$ such that
\begin{equation}\label{eq:f-hyp}
f(t)\ge2\quad\text{and}\quad f'(t)\ge1\qquad\text{for all real }t\ge1 .
\end{equation}
In particular $f$ has positive leading coefficient and is strictly increasing on $[1,\infty)$. We study
\begin{equation}\label{eq:Vf}
V_f(x,y)=\sum_{1\le k\le y}f'(k)\,V(x;f(k))=\sum_{1\le k\le y}f'(k)\sum_{l=1}^{f(k)}E(x;f(k),l)^2,\qquad Y:=f(y)\le x,
\end{equation}
where $k$ runs over integers and $y\ge1$ is real. The weight $f'(k)$ compensates for the sparseness of the moduli: by the mean value theorem, the number of integers in an interval $(f(k),f(k+1)]$ is about $f'(k)$, so $\sum_{k\le y}f'(k)G(f(k))$ is a Riemann-sum approximation to $\int_{f(1)}^{Y}G(u)\,du$, and $V_f(x,y)$ is directly comparable with the full variance $V(x,Y)$.

The question is to determine the asymptotic behaviour of $V_f(x,y)$, with an explicit main term depending on $f$, in as large a range of $f(y)\le x$ as possible.

\subsection{Background}

For the full variance, Vaughan~\cite[Theorem, (1.12)--(1.15)]{Vaughan2005} proved, for $k$-free numbers in general and in the case $k=2$ of squarefree numbers in particular, that whenever $Y\le x$,
\[
V(x,Y)=c_2\,x^{1/2}Y^{3/2}+O\Bigl(x^{1/4}Y^{7/4}\exp\Bigl(-c^*\tfrac{(\log 2x/Y)^{3/5}}{(\log\log 3x/Y)^{1/5}}\Bigr)\Bigr)+O\Bigl(x^{3/2}\exp\Bigl(-\tilde c\tfrac{(\log x)^{3/5}}{(\log\log x)^{1/5}}\Bigr)\Bigr),
\]
with
\begin{equation}\label{eq:vaughan-c2}
c_2=\frac{8C_1\,(-\zeta(-1/2))}{3\,\zeta(2)^2},\qquad
C_1=\prod_p\Bigl(1+\frac{(p-1)^2p^{3/2}-2p^4+3p^3-p}{p^2(p^2-1)^2}\Bigr).
\end{equation}
(This is (1.14)--(1.15) of~\cite{Vaughan2005} with $k=2$, written with $Y$ in place of Vaughan's modulus bound $Q$, to avoid a clash with the counting function $Q(x)$; there $x$ is assumed sufficiently large.) Vaughan's main term dominates when $x/Y$ is large and $Y$ is large compared with $x^{2/3}\exp\bigl(-\tfrac23\tilde c(\log x)^{3/5}(\log\log x)^{-1/5}\bigr)$, a quantity slightly smaller than $x^{2/3}$~\cite[(1.16)]{Vaughan2005}. His proof uses the circle method, with major-arc information on the exponential sum over $k$-free numbers from~\cite{BGPVW1998}. Earlier results, quoted in~\cite[(1.6)--(1.11)]{Vaughan2005}, have the shape $V(x,Y)=x^2\mathcal F(x/Y)+O(\cdots)$, with an error term of size at least about $x^{3/2}$.

For primes, Brüdern and Wooley~\cite{BruedernWooley2011} proved an analogue of the Montgomery--Hooley formula in which the moduli are restricted to the values $f(k)$ of a polynomial of degree at least $2$, with the same weight $f'(k)$. They adapted the circle-method approach to the Montgomery--Hooley formula that they attribute to Goldston and Vaughan~\cite{GoldstonVaughan1997} and to Vaughan~\cite{Vaughan1998I}. The problem considered here is the squarefree analogue of their result.

The argument presented below uses no circle method. It is completely elementary, and the price is a narrower range of validity of the asymptotic formula, namely $x^{7/9+\delta}\le f(y)\le x^{1-\delta}$; see \cref{rem:range}.

\subsection{The main result}

We need some notation. For a prime $p$ let
\begin{equation}\label{eq:abc}
a_p=1-\frac3{p^2}+\frac2{p^3}=\frac{(p-1)^2(p+2)}{p^3},\qquad
b_p=\frac{(p-1)(p^2-2+\sqrt p)}{p^3},\qquad
c_p=\frac{(p-1)^2(p+1)}{p^3}.
\end{equation}
Let $\rho(m)=\#\{a\bmod m:\ f(a)\equiv0\ (\mathrm{mod}\ m)\}$, and put
\begin{equation}\label{eq:sigma}
\sigma_p=a_p\Bigl(1-\frac{\rho(p)}p\Bigr)+b_p\Bigl(\frac{\rho(p)}p-\frac{\rho(p^2)}{p^2}\Bigr)+c_p\,\frac{\rho(p^2)}{p^2},
\qquad
C_f=\frac23\,I\prod_p\sigma_p,
\end{equation}
where
\begin{equation}\label{eq:I}
I=\int_0^\infty B(u^{-2})\,du=-4\zeta(-\tfrac12)=\frac{\zeta(3/2)}{\pi}=\IConst\ldots,\qquad B(u)=\fr{u}(1-\fr{u}),
\end{equation}
and $\fr u$ is the fractional part of $u$.
The three weights in $\sigma_p$ are the proportions of residues $k$ modulo $p^2$ with $p\nmid f(k)$, with $p\,\|\,f(k)$ and with $p^2\mid f(k)$, respectively.
We also need the explicit single-modulus main term
\begin{equation}\label{eq:MT}
\MT(x;q)=A_0\,q\sum_{t\ge1}\eta(t)\,m_t(q)\,B\Bigl(\frac{x}{q\,m_t(q)}\Bigr),
\end{equation}
where
\begin{equation}\label{eq:A0-eta}
A_0=\prod_p\Bigl(1-\frac2{p^2}\Bigr),\qquad \eta(t)=\mu^2(t)\prod_{p\mid t}\frac1{p^2-2},\qquad m_t(q)=\frac{t^2}{\gcd(t^2,q)}.
\end{equation}
The series~\eqref{eq:MT} converges absolutely, since $0\le B(u)\le u$ (\cref{lem:eta}).

\begin{theorem}\label{thm:main}
Let $f\in\Z[t]$ have degree $d\ge1$ and satisfy~\eqref{eq:f-hyp}. The product $\prod_p\sigma_p$ converges absolutely, and $0<C_f<\infty$. For every $\eps>0$ there is a constant $c(f,\eps)$, depending only on $f$ and $\eps$, such that for all real $x\ge2$ and $y\ge1$ with $Y=f(y)\le x$:
\begin{enumerate}[label=\textup{(\roman*)}]
\item $\displaystyle\bigl|V_f(x,y)-C_f\,x^{1/2}Y^{3/2}\bigr|\le c(f,\eps)\Bigl(x^{1/3+\eps}Y^{5/3}+x^{5/3+\eps}+x^{1/2+\eps}Y^{3/2-1/(2d)}\Bigr);$
\item $\displaystyle\Bigl|V_f(x,y)-\sum_{1\le k\le y}f'(k)\,\MT(x;f(k))\Bigr|\le c(f,\eps)\bigl(x^{5/3+\eps}+x^{1/2}Y\bigr).$
\end{enumerate}
\begin{enumerate}[label=\textup{(\roman*)},resume]
\item Consequently, for each fixed $\delta\in(0,1/9)$, uniformly for $x^{7/9+\delta}\le f(y)\le x^{1-\delta}$,
\[
V_f(x,y)=C_f\,x^{1/2}f(y)^{3/2}\bigl(1+O_{f,\delta}(x^{-\eta})\bigr),\qquad \eta=\min\Bigl(\frac{\delta}{12},\frac1{6d}\Bigr).
\]
\end{enumerate}
\end{theorem}

\begin{remark}[The range]\label{rem:range}
Each error term in (i) comes from one step of the proof. The term $x^{5/3+\eps}$ comes from the pointwise bound $x^{2/3+\eps}$ for the error in the count of pairs of squarefree numbers with a given difference (\cref{prop:pairs}), summed trivially over about $x/q$ differences. It is smaller than $x^{1/2}Y^{3/2}$ only when $Y\ge x^{7/9+\delta}$. The term $x^{1/3+\eps}Y^{5/3}$ comes from the Euler--Maclaurin error $Z^{1/3}$ in the sawtooth sums of \cref{prop:sawtooth}. It has relative size $x^\eps(Y/x)^{1/6}$, which forces $Y\le x^{1-\delta}$. The term $x^{1/2+\eps}Y^{3/2-1/(2d)}$ comes from averaging over the polynomial values (\cref{prop:poly}). Statement~(ii) avoids the sawtooth asymptotics altogether: its error term is of smaller order than $x^{1/2}Y^{3/2}$ throughout $x^{7/9+\delta}\le Y\le x$, including $Y\asymp x$, where the arithmetic main term $\sum_kf'(k)\MT(x;f(k))$ is kept as it is.
\end{remark}

\begin{remark}[Consistency with the full variance]\label{rem:vaughan}
Formally taking $f(t)=t$ in~\eqref{eq:sigma}, so that $\rho\equiv1$, gives $\sigma_p=C_{1,p}(1-p^{-2})^2$ identically in $p$, where $C_{1,p}$ is the factor of $C_1$ in~\eqref{eq:vaughan-c2}. Since $\prod_p(1-p^{-2})^2=\zeta(2)^{-2}$, the constant $\frac23I\prod_p\sigma_p$ equals Vaughan's $c_2$ exactly. The identity of local factors was checked with computer algebra, and both products are evaluated numerically in \cref{tab:constants}. The polynomial $f(t)=t$ violates $f(1)\ge2$, so this is a consistency check and not a case of \cref{thm:main}.
\end{remark}

\subsection{Strategy}

The proof has four parts.

\emph{An exact identity for one modulus} (\cref{sec:identity}). Expanding the square, $\sum_lQ(x;q,l)^2$ counts pairs $n\equiv n'\pmod q$ of squarefree numbers up to $x$, that is, the diagonal together with pairs $(n,n+hq)$. Replace the number of pairs with difference $m$ by its expected value $A(m)(x-m)$, where $A(m)$ is the density of pairs of squarefree numbers with difference $m$. Summing $A(hq)(x-hq)$ over $h$ is an arithmetic-progression sum. The terms of order $x^2/q$ and $x$ in this sum cancel \emph{exactly} against $x^2\sum_lg(q,l)^2$ and the diagonal. The remainder is the lattice-point discrepancy $\MT(x;q)$, which arises from the identity $2\sum_{1\le j\le T}(T-j)=T^2-T+B(T)$. What is left is controlled by three explicit error terms.

\emph{Error terms} (\cref{sec:densities,sec:pairs}). The pair counts are evaluated with a uniform error $O(x^{2/3+\eps})$ by a Möbius truncation at $x^{1/3}$ (\cref{prop:pairs}). The remaining terms are of size $O(x^{1/2})$ and $O(x^{3/2}/q)$ (\cref{prop:trivial}).

\emph{The main term} (\cref{sec:sawtooth}). After a factorisation of $t$ according to $q$, $\MT(x;q)$ is a weighted sum of sawtooth sums $\sum_sB(Z/s^2)$. By Euler--Maclaurin summation such a sum equals $I\sqrt Z+O(Z^{1/3})$, which gives $\MT(x;q)=h(q)\sqrt{xq}+O_\eps(x^{1/3}q^{2/3+\eps})$ with an explicit multiplicative-type function $h$ (\cref{prop:MT}). This is the source of the order $x^{1/2}q^{1/2}$ of $V(x;q)$.

\emph{Averaging over polynomial values} (\cref{sec:poly}). Finally, $\sum_{k\le y}f'(k)f(k)^{1/2}h(f(k))$ is evaluated by expanding $h$ as a divisor sum and counting $k$ in the $\rho(e)$ residue classes modulo $e$ on which $e\mid f(k)$. This gives $C_fY^{3/2}$, where $\tfrac23Y^{3/2}$ comes from $\int f'f^{1/2}=\tfrac23f^{3/2}$.

\Cref{sec:assembly} combines these steps and composes the error terms explicitly. \Cref{sec:numerics} compares \cref{thm:main} with exact computations of $V_f(x,y)$ for three polynomials.

\subsection{Notation and conventions}\label{sec:notation}

The letter $p$ always denotes a prime. We write $p^j\,\|\,q$ if $p^j\mid q$ and $p^{j+1}\nmid q$, $v_p(q)$ for the exponent of $p$ in $q$, $\omega(n)$ for the number of distinct prime factors of $n$, $\tau(n)$ for the number of divisors of $n$, and $\rad(n)=\prod_{p\mid n}p$. Further, $\one[\cdot]$ is the indicator of a statement, $(a,b)=\gcd(a,b)$ and $[a,b]=\lcm(a,b)$, and $\fr u=u-\lfloor u\rfloor$. A quantity written $\theta$ (possibly with an index) is a real number with $|\theta|\le1$, which may differ at each occurrence. Implied constants in $O$, $\ll$ and $\gg$ and constants $c(\cdot)$ depend only on the parameters shown. The symbols $B$, $I$, $A_0$, $\eta$, $m_t(q)$, $\MT$, $a_p$, $b_p$, $c_p$, $\rho$, $\sigma_p$ and $C_f$ have the fixed meanings given above. In addition,
\begin{equation}\label{eq:kappa-S2}
\kappa(t)=t^2\eta(t)=\mu^2(t)\prod_{p\mid t}\frac{p^2}{p^2-2},\qquad S_2(q)=\sum_{l=1}^qg(q,l)^2,\qquad F(t)=f'(t)f(t)^{1/2},
\end{equation}
and the multiplicative-type function
\begin{equation}\label{eq:h}
h(q)=I\prod_{p\nmid q}a_p\prod_{p\,\|\,q}b_p\prod_{p^2\mid q}c_p .
\end{equation}

\section{Preliminaries}\label{sec:prelim}

This section collects the elementary tools used repeatedly later. The only external inputs are standard facts from Montgomery and Vaughan~\cite{MV2007}, cited precisely where they are used.

\subsection{Counting}

\begin{lemma}[A tail bound]\label{lem:tail}
For every real $w>0$, $\displaystyle\sum_{a>w}a^{-2}\le\frac2w$.
\end{lemma}

\begin{proof}
If $w\ge1$, let $k=\lfloor w\rfloor+1>w$. Then $\sum_{a\ge k}a^{-2}\le k^{-2}+\int_k^\infty u^{-2}\,du\le 2/k<2/w$. If $0<w<1$, the sum is at most $\zeta(2)<2<2/w$.
\end{proof}

\begin{lemma}[Integers in a residue class]\label{lem:class}
Let $N\ge0$ be real, $L\in\N$ and $c\in\Z$. Then $\#\{1\le n\le N:\ n\equiv c\ (\mathrm{mod}\ L)\}=N/L+\theta$ with $|\theta|<1$.
\end{lemma}

\begin{proof}
Let $n_0\in\{1,\dots,L\}$ be the least positive element of the class. If $n_0>N$, the count is $0$, and $0\le N/L<n_0/L\le1$. Otherwise the count is $\lfloor(N-n_0)/L\rfloor+1$, which lies in $\bigl((N-n_0)/L,(N-n_0)/L+1\bigr]$. This interval is contained in $(N/L-1,N/L+1)$, because $1\le n_0\le L$.
\end{proof}

\begin{lemma}[Divisor bounds]\label{lem:divisor}
For every $\eps>0$ there is $C(\eps)$ with $\tau(n)\le C(\eps)n^\eps$ for all $n\in\N$~\cite[(2.20)]{MV2007}. Moreover $2^{\omega(n)}\le\tau(n)$ and $\sum_{e\mid n}4^{\omega(e)}\le\tau(n)^3$.
\end{lemma}

\begin{proof}
The first statement is~\cite[(2.20)]{MV2007}. Each subset of the prime divisors of $n$ gives a divisor, so $2^{\omega(n)}\le\tau(n)$. For $e\mid n$ we have $4^{\omega(e)}=(2^{\omega(e)})^2\le\tau(e)^2\le\tau(n)^2$, and there are $\tau(n)$ divisors $e$.
\end{proof}

\subsection{Euler products}

We use the following standard fact in the form given by Montgomery and Vaughan~\cite[Theorem~1.9]{MV2007} (with $s=0$): \emph{if $G$ is multiplicative and $\sum_{n\ge1}|G(n)|<\infty$, then}
\begin{equation}\label{eq:euler}
\sum_{n\ge1}G(n)=\prod_p\Bigl(1+\sum_{j\ge1}G(p^j)\Bigr).
\end{equation}
We also use that if $\sum_p|u_p|<\infty$ and $1+u_p\ne0$ for all $p$, then $\prod_p(1+u_p)$ converges absolutely to a nonzero limit, and such products may be multiplied and divided factor by factor. Absolute convergence of $\sum|G(n)|$ for a multiplicative $G$ follows from $\prod_p\bigl(1+\sum_j|G(p^j)|\bigr)<\infty$, because every partial sum of $\sum|G(n)|$ is bounded by a finite subproduct.

\begin{lemma}[The functions $\eta$ and $\kappa$]\label{lem:eta}
The functions $\eta$ and $\kappa$ of~\eqref{eq:A0-eta} and~\eqref{eq:kappa-S2} are multiplicative, nonnegative, and supported on squarefree integers. For all $t$,
\begin{equation}\label{eq:eta-bound}
0\le\eta(t)\le2^{\omega(t)}t^{-2},\qquad 0\le\kappa(t)\le2^{\omega(t)}.
\end{equation}
Further, $0<A_0\le1$ and
\begin{equation}\label{eq:sum-eta}
\sum_{t\ge1}\eta(t)=\prod_p\Bigl(1+\frac1{p^2-2}\Bigr)=\frac{6}{\pi^2A_0}.
\end{equation}
For all real $x>0$ and $q\in\N$, the terms of the series~\eqref{eq:MT} satisfy $0\le\eta(t)m_t(q)B\bigl(x/(qm_t(q))\bigr)\le\eta(t)\,x/q$. Hence $\MT(x;q)$ converges absolutely and $0\le\MT(x;q)\le 6x/\pi^2$.
\end{lemma}

\begin{proof}
Multiplicativity and the support are clear from the definitions. Since $p^2-2\ge p^2/2$ for every prime, $\eta(p)\le2/p^2$ and $\kappa(p)\le2$, which gives~\eqref{eq:eta-bound}. Each factor $1-2/p^2$ lies in $[\tfrac12,1)$ and $\sum_p2/p^2<\infty$, so $0<A_0\le1$. By~\eqref{eq:euler}, applied to $\eta$ (whose series converges since $\sum_t2^{\omega(t)}t^{-2}\le\sum_t\tau(t)t^{-2}=\zeta(2)^2$ by~\cite[Theorem~1.8]{MV2007}),
\[
\sum_t\eta(t)=\prod_p\Bigl(1+\frac1{p^2-2}\Bigr)=\prod_p\frac{1-p^{-2}}{1-2p^{-2}}=\frac{\prod_p(1-p^{-2})}{A_0}=\frac{6}{\pi^2A_0},
\]
using $\prod_p(1-p^{-2})=1/\zeta(2)=6/\pi^2$. Finally $0\le B(u)\le\min(u,\frac14)$ for $u\ge0$, so $m\,B(x/(qm))\le x/q$, and $\MT(x;q)\le A_0q\cdot(x/q)\sum_t\eta(t)=6x/\pi^2$.
\end{proof}

The next evaluation occurs twice, in the second moment of $g$ (\cref{prop:densities}(iv)) and in the density of pairs of squarefree numbers (\cref{prop:pairs}).

\begin{lemma}[Coprime pairs of squarefree numbers]\label{lem:pairs-euler}
For every $e\in\N$,
\[
\sum_{\substack{a,b\ge1,\ \mu^2(ab)=1\\ (ab,e)=1}}\frac{\mu(a)\mu(b)}{a^2b^2}
=\prod_{p\nmid e}\Bigl(1-\frac2{p^2}\Bigr)=A_0\prod_{p\mid e}\frac{p^2}{p^2-2},
\]
and the series converges absolutely. Here $\mu^2(ab)=1$ means exactly that $a$ and $b$ are squarefree and coprime.
\end{lemma}

\begin{proof}
Group the pairs according to $n=ab$. For a squarefree $n$ there are $2^{\omega(n)}$ pairs $(a,b)$ with $ab=n$, each with $\mu(a)\mu(b)=\mu(n)$. The sum is therefore $\sum_{n}G(n)$ with $G(n)=\mu(n)2^{\omega(n)}n^{-2}\one[(n,e)=1]$, which is multiplicative with $\sum|G(n)|\le\zeta(2)^2$. By~\eqref{eq:euler} the sum equals $\prod_{p\nmid e}(1-2/p^2)$. The last expression follows by dividing $A_0$ by the finitely many factors $1-2/p^2=(p^2-2)/p^2\ne0$ with $p\mid e$.
\end{proof}

\subsection{Euler--Maclaurin summation}

Put $\psi(u)=\fr u-\frac12$ and
\begin{equation}\label{eq:psi2}
\psi_2(u)=\tfrac12\bigl(\fr u^2-\fr u\bigr)+\tfrac1{12}=\tfrac1{12}-\tfrac12B(u).
\end{equation}
Then $\psi_2$ is continuous and $1$-periodic, $\psi_2'=\psi$ off the integers, $|\psi_2|\le\frac1{12}$ and $|\psi|\le\frac12$, and $\psi_2(n)=\frac1{12}$, $\psi(n)=-\frac12$ for $n\in\Z$.

\begin{lemma}[Euler--Maclaurin]\label{lem:EM}
Let $a<b$ be real and let $G\in C^2[a,b]$ be real- or complex-valued. Then
\[
\sum_{a<n\le b}G(n)=\int_a^bG(u)\,du+\psi(a)G(a)-\psi(b)G(b)-\psi_2(a)G'(a)+\psi_2(b)G'(b)-\int_a^b\psi_2(u)G''(u)\,du .
\]
Suppose now that $G$ is $C^2$ on $[a,\infty)$, that $G(u)\to0$ and $G'(u)\to0$ as $u\to\infty$, and that $G''\in L^1[a,\infty)$. Then
\[
\lim_{b\to\infty}\Bigl(\sum_{a<n\le b}G(n)-\int_a^bG(u)\,du\Bigr)=\psi(a)G(a)-\psi_2(a)G'(a)-\int_a^\infty\psi_2(u)G''(u)\,du .
\]
In particular, if one of $\sum_{n>a}G(n)$ and $\int_a^\infty G(u)\,du$ converges, then so does the other, and the formula holds with $b=\infty$ and the terms at $b$ omitted.
\end{lemma}

\begin{proof}
Since $\lfloor u\rfloor=u-\frac12-\psi(u)$, we have $\sum_{a<n\le b}G(n)=\int_{(a,b]}G\,d\lfloor u\rfloor=\int_a^bG\,du-\int_{(a,b]}G\,d\psi$. As $\psi$ is right-continuous and of bounded variation and $G$ is continuous, integration by parts for Riemann--Stieltjes integrals gives $\int_{(a,b]}G\,d\psi=G(b)\psi(b)-G(a)\psi(a)-\int_a^b\psi G'\,du$. Finally $\int_a^b\psi G'\,du=[\psi_2G']_a^b-\int_a^b\psi_2G''\,du$. (As a check, $G\equiv1$ gives $\lfloor b\rfloor-\lfloor a\rfloor$.) For complex-valued $G$, apply this to the real and imaginary parts.

For $b\to\infty$, move $\int_a^bG\,du$ to the left-hand side of the finite formula. On the right, the terms at $b$ tend to $0$, because $\psi$ and $\psi_2$ are bounded and $G(b),G'(b)\to0$. The last integral converges to $\int_a^\infty\psi_2G''\,du$, because $\psi_2$ is bounded and $G''\in L^1$. This proves the limit formula. The partial sums $\sum_{a<n\le b}G(n)$, for real $b$, are those of the series $\sum_{n>a}G(n)$. So if the series converges, then $\int_a^bG\,du$ converges as $b\to\infty$, and conversely; in either case the limit formula becomes the stated formula with $b=\infty$. (The series and the integral need not converge separately: $G(u)=\cos(\log u)/u$ on $[1,\infty)$ satisfies the hypotheses, but $\int_1^XG\,du=\sin(\log X)$ has no limit as $X\to\infty$, and hence neither have the partial sums of $\sum_{n\ge2}G(n)$.)
\end{proof}

In the applications below (\cref{prop:sawtooth}), $G$ is $u^{-s}$ with $\Re s>1$ or a linear combination of $u^{-2}$ and $u^{-4}$. In both cases the series and the integral converge absolutely, so the last statement of the lemma applies.

\Cref{lem:EM} is the Euler--Maclaurin formula of the first and second order; compare~\cite[§1.3]{MV2007}, where the continuation of $\zeta(s)$ to $\Re s>0$ is obtained in the same way~\cite[Theorem~1.12]{MV2007}.

\section{Densities of squarefree numbers in progressions}\label{sec:densities}

We first determine $g(q,l)$ and its first two moments over $l$. We then show that two of the four terms in the decomposition of \cref{sec:identity} are small. Throughout this section $q\in\N$ is fixed, and for $d\in\N$ we write
\[
g_d=\gcd(d^2,q).
\]

\begin{proposition}[Closed form and moments of $g$]\label{prop:densities}
For every $q\in\N$ and $l\in\Z$:
\begin{enumerate}[label=\textup{(\roman*)}]
\item The limit $g(q,l)=\lim_{x\to\infty}Q(x;q,l)/x$ exists, and
\begin{equation}\label{eq:g-series}
g(q,l)=\sum_{d\ge1}\mu(d)\,\frac{g_d}{d^2q}\,\one[g_d\mid l],
\end{equation}
where the series converges absolutely. Moreover $0\le g(q,l)\le1/q$, and $g(q,l)$ depends only on $l\bmod q$.
\item $\sum_{l=1}^qg(q,l)=6/\pi^2$.
\item For every squarefree $n\in\N$,
\[
g(q,n)=\frac{6}{\pi^2}\,\frac1q\prod_{p\mid q}\Bigl(1-\frac1{p^2}\Bigr)^{-1}\prod_{\substack{p\mid(n,q)\\ p\,\|\,q}}\Bigl(1-\frac1p\Bigr).
\]
\item $S_2(q)=\sum_{l=1}^qg(q,l)^2=\dfrac{A_0}{q}\displaystyle\sum_{t\ge1}\eta(t)\frac{\gcd(t^2,q)}{t^2}=\frac{A_0}{q}\sum_{t\ge1}\frac{\eta(t)}{m_t(q)}$.
\item $\sum_{t\ge1}\eta(t)=6/(\pi^2A_0)$.
\end{enumerate}
\end{proposition}

\begin{proof}
(i) We use $\mu^2(n)=\sum_{d^2\mid n}\mu(d)$. For $d>\sqrt x$ no $n\le x$ is divisible by $d^2$, so
\[
Q(x;q,l)=\sum_{d\le\sqrt x}\mu(d)N_d(x),\qquad N_d(x)=\#\{n\le x:\ n\equiv l\ (q),\ d^2\mid n\}.
\]
Write $n=d^2m$ with $1\le m\le x/d^2$. The congruence $d^2m\equiv l\pmod q$ is solvable in $m$ if and only if $g_d\mid l$. If $g_d\nmid l$, then $N_d(x)=0$. If $g_d\mid l$, the congruence is equivalent to $(d^2/g_d)m\equiv l/g_d\pmod{q/g_d}$. As $(d^2/g_d,q/g_d)=1$, this is a single residue class modulo $q/g_d$, and \cref{lem:class} gives $N_d(x)=\frac{g_d}{d^2q}x+\theta_d$. Hence
\[
Q(x;q,l)=x\sum_{d\le\sqrt x}\mu(d)\frac{g_d}{d^2q}\one[g_d\mid l]+\theta\sqrt x .
\]
Since $g_d\le q$, the $d$-th term of~\eqref{eq:g-series} has modulus at most $d^{-2}$. So the series converges absolutely, and by \cref{lem:tail} its tail beyond $\sqrt x$ is at most $2/\sqrt x$ for $x\ge1$. Therefore
\begin{equation}\label{eq:Q-g}
|Q(x;q,l)-x\,g(q,l)|\le3\sqrt x\qquad(x\ge1),
\end{equation}
with $g(q,l)$ defined by~\eqref{eq:g-series}, and $Q(x;q,l)/x\to g(q,l)$. From $0\le Q(x;q,l)\le x/q+1$ we get $0\le g(q,l)\le1/q$. Since $g_d\mid q$, the condition $g_d\mid l$ depends only on $l\bmod q$.

(ii) For each $d$, exactly $q/g_d$ of the integers $1\le l\le q$ are divisible by $g_d$. Summing the absolutely convergent series~\eqref{eq:g-series} over the $q$ values of $l$ and interchanging the sums gives $\sum_l g(q,l)=\sum_d\mu(d)d^{-2}=6/\pi^2$.

(iii) Let $n$ be squarefree. Only squarefree $d$ contribute to~\eqref{eq:g-series}. For squarefree $d$, $g_d=\prod_{p\mid d}\gcd(p^2,q)$ is a product of coprime prime powers, so $g_d\mid n$ if and only if $\gcd(p^2,q)\mid n$ for every $p\mid d$. Hence, for squarefree $d$,
\[
\mu(d)\frac{g_d}{d^2}\one[g_d\mid n]=\prod_{p\mid d}u(p),\qquad u(p)=-\frac{\gcd(p^2,q)}{p^2}\one[\gcd(p^2,q)\mid n],
\]
and $q\,g(q,n)=\sum_{d}\mu^2(d)\prod_{p\mid d}u(p)$. We have $|u(p)|=p^{-2}$ for $p\nmid q$, and $|u(p)|\le1$ for the finitely many $p\mid q$. So $\sum_p|u(p)|<\infty$ and, by~\eqref{eq:euler}, $q\,g(q,n)=\prod_p(1+u(p))$. Now
\begin{itemize}
\item if $p\nmid q$, then $\gcd(p^2,q)=1$ and $1+u(p)=1-p^{-2}$;
\item if $p\,\|\,q$, then $\gcd(p^2,q)=p$ and $1+u(p)=1-p^{-1}\one[p\mid n]$;
\item if $p^2\mid q$, then $\gcd(p^2,q)=p^2\nmid n$, since $n$ is squarefree, and $1+u(p)=1$.
\end{itemize}
Hence $q\,g(q,n)=\prod_{p\nmid q}(1-p^{-2})\prod_{p\|q,\ p\mid n}(1-p^{-1})$. The claim follows from $\prod_{p\nmid q}(1-p^{-2})=\frac6{\pi^2}\prod_{p\mid q}(1-p^{-2})^{-1}$.

(iv) Squaring~\eqref{eq:g-series} (a product of two absolutely convergent series) and summing over $1\le l\le q$ gives
\[
S_2(q)=\frac1{q^2}\sum_{d,e\ge1}\mu(d)\mu(e)\frac{g_dg_e}{d^2e^2}\,\#\{1\le l\le q:\ g_d\mid l,\ g_e\mid l\}.
\]
Both $g_d$ and $g_e$ divide $q$, so the count is $q/[g_d,g_e]$, and $g_dg_e/[g_d,g_e]=(g_d,g_e)=\gcd\bigl((d,e)^2,q\bigr)$. Thus, with $G(t)=\gcd(t^2,q)\le q$,
\[
S_2(q)=\frac1q\sum_{d,e\ \mathrm{squarefree}}\mu(d)\mu(e)\frac{G((d,e))}{d^2e^2},
\]
which converges absolutely. Write $t=(d,e)$, $d=td'$ and $e=te'$. This is a bijection from pairs of squarefree $d,e$ onto triples $(t,d',e')$ with $t$ squarefree, $d'e'$ squarefree and $(d'e',t)=1$, and $\mu(d)\mu(e)=\mu(d')\mu(e')$. Summing over $t$ last and using \cref{lem:pairs-euler} with $e=t$,
\[
S_2(q)=\frac1q\sum_{t\ \mathrm{squarefree}}\frac{G(t)}{t^4}A_0\prod_{p\mid t}\frac{p^2}{p^2-2}=\frac{A_0}{q}\sum_{t\ge1}\eta(t)\frac{\gcd(t^2,q)}{t^2},
\]
because $t^{-4}\prod_{p\mid t}p^2/(p^2-2)=t^{-2}\eta(t)$ for squarefree $t$. Finally $\gcd(t^2,q)/t^2=1/m_t(q)$.

(v) This is~\eqref{eq:sum-eta}.
\end{proof}

For $q=1$, (iv) reads $A_0\sum_t\eta(t)t^{-2}=\prod_p(1-p^{-2})^2=36/\pi^4=g(1,1)^2$, as it should.

The second ingredient is an elementary count of squarefree numbers coprime to a given squarefree number.

\begin{lemma}\label{lem:Tr}
Let $r$ be squarefree and put $T_r(y)=\#\{1\le m\le y:\ \mu^2(m)=1,\ (m,r)=1\}$ for real $y>0$. Then
\[
|T_r(y)-\nu_ry|\le3\cdot2^{\omega(r)}\sqrt y,\qquad \nu_r=\sum_{d\ge1}\sum_{e\mid r}\frac{\mu(d)\mu(e)}{[d^2,e]},
\]
where the double series converges absolutely and $|\nu_r|\le2\cdot2^{\omega(r)}$. In particular $\nu_1=6/\pi^2$, so $|Q(y)-6y/\pi^2|\le3\sqrt y$ for all $y>0$.
\end{lemma}

\begin{proof}
Insert $\mu^2(m)=\sum_{d^2\mid m}\mu(d)$ and $\one[(m,r)=1]=\sum_{e\mid(m,r)}\mu(e)$. This gives $T_r(y)=\sum_{d\ge1}\sum_{e\mid r}\mu(d)\mu(e)\lfloor y/[d^2,e]\rfloor$, where only $d\le\sqrt y$ contribute because $[d^2,e]\ge d^2$. The same inequality gives $|\nu_r|\le2^{\omega(r)}\zeta(2)\le2\cdot2^{\omega(r)}$. If $y\ge1$, there are at most $2^{\omega(r)}\sqrt y$ pairs $(d,e)$ with $d\le\sqrt y$, $e\mid r$. Replacing each floor by $y/[d^2,e]$ costs at most $1$ per pair. Completing the $d$-sum costs at most $y\,2^{\omega(r)}\sum_{d>\sqrt y}d^{-2}\le2\cdot2^{\omega(r)}\sqrt y$, by \cref{lem:tail}. If $0<y<1$, then $T_r(y)=0$ and $|\nu_ry|\le2\cdot2^{\omega(r)}y\le2\cdot2^{\omega(r)}\sqrt y$. For $r=1$ only $e=1$ occurs, and $\nu_1=\sum_d\mu(d)d^{-2}=6/\pi^2$.
\end{proof}

\begin{proposition}[Two small terms]\label{prop:trivial}
Let $c_0=3e^{2\zeta(3/2)}$. For all real $x\ge1$ and all integers $1\le q\le x$:
\begin{enumerate}[label=\textup{(\roman*)}]
\item $|Q(x)-6x/\pi^2|\le c_0\sqrt x$;
\item $\Bigl|\sum_{l=1}^qg(q,l)E(x;q,l)\Bigr|\le c_0\sqrt x/q$.
\end{enumerate}
\end{proposition}

The hypothesis $q\le x$ is not used in the proof of (ii).

\begin{proof}
(i) This is the case $r=1$ of \cref{lem:Tr}, since $3\le c_0$.

(ii) Let $q_1=\prod_{p\|q}p$. Since $g(q,\cdot)$ is $q$-periodic,
\begin{equation}\label{eq:gE}
S(x):=\sum_{l=1}^qg(q,l)Q(x;q,l)=\sum_{n\le x}\mu^2(n)g(q,n),\qquad \sum_{l=1}^qg(q,l)E(x;q,l)=S(x)-x\,S_2(q).
\end{equation}
For squarefree $n$ the last product in \cref{prop:densities}(iii) runs over $p\mid(n,q_1)$, and $\prod_{p\mid(n,q_1)}(1-1/p)=\sum_{r\mid(n,q_1)}\mu(r)/r$. Put $c_q=\frac6{\pi^2}\prod_{p\mid q}(1-p^{-2})^{-1}$; then $0<c_q\le\frac6{\pi^2}\zeta(2)=1$ and $g(q,n)=\frac{c_q}{q}\sum_{r\mid q_1,\ r\mid n}\mu(r)/r$. Interchanging the finite sum over $r$ with the sum over $n$,
\[
S(x)=\frac{c_q}{q}\sum_{r\mid q_1}\frac{\mu(r)}{r}N_r(x),\qquad N_r(x)=\#\{n\le x:\ r\mid n,\ \mu^2(n)=1\}.
\]
For squarefree $r$, $n=rm$ is squarefree if and only if $m$ is squarefree and $(m,r)=1$, so $N_r(x)=T_r(x/r)$. By \cref{lem:Tr},
\[
S(x)=M_q\,x+R_q(x),\qquad M_q=\frac{c_q}{q}\sum_{r\mid q_1}\frac{\mu(r)\nu_r}{r^2},
\qquad
|R_q(x)|\le\frac{3\sqrt x}{q}\sum_{r\ \mathrm{squarefree}}\frac{2^{\omega(r)}}{r^{3/2}}.
\]
The last sum equals $\prod_p(1+2p^{-3/2})\le\exp\bigl(2\sum_pp^{-3/2}\bigr)\le e^{2\zeta(3/2)}$, so $|R_q(x)|\le c_0\sqrt x/q$ for all $x>0$.

It remains to show $M_q=S_2(q)$. As $x\to\infty$, $Q(x;q,l)/x\to g(q,l)$ for each of the finitely many $l$, so $S(x)/x\to S_2(q)$. On the other hand $S(x)/x=M_q+O(x^{-1/2})$. Hence $M_q=S_2(q)$, and by~\eqref{eq:gE} the sum $\sum_lg(q,l)E(x;q,l)$ equals $R_q(x)$.
\end{proof}

\section{An exact decomposition of the single-modulus variance}\label{sec:identity}

For real $x\ge1$ and an integer $m\ge1$ let
\begin{gather}\label{eq:pairs-def}
C_x(m)=\#\{n\ge1:\ n+m\le x,\ \mu^2(n)=\mu^2(n+m)=1\},\\
A(m)=A_0\sum_{t^2\mid m}\eta(t),\qquad D_x(m)=C_x(m)-A(m)(x-m).\notag
\end{gather}
Here $A(m)$ is the density of integers $n$ with $n$ and $n+m$ both squarefree (\cref{prop:pairs}), so $A(m)(x-m)$ is the expected value of $C_x(m)$ and $D_x(m)$ is its error.

\begin{proposition}[Exact decomposition]\label{prop:identity}
For every real $x\ge1$ and every $q\in\N$,
\[
V(x;q)=\MT(x;q)+2\sum_{1\le h\le x/q}D_x(hq)+\Bigl(Q(x)-\frac{6x}{\pi^2}\Bigr)-2x\sum_{l=1}^qg(q,l)E(x;q,l).
\]
The sum over $h$ is empty if $q>x$.
\end{proposition}

All four terms have a transparent meaning. The second term collects the errors in the counts of pairs, and the last two are of size $O(\sqrt x)$ and $O(x^{3/2}/q)$ by \cref{prop:trivial}. The main term $\MT(x;q)$ is what remains of the \emph{expected} pair counts after an exact cancellation of the terms of order $x^2/q$ and $x$.

\begin{proof}
Write $Q_l=Q(x;q,l)$, $g_l=g(q,l)$ and $E_l=Q_l-g_lx$. Expanding the square and using $\sum_lg_lQ_l=\sum_lg_lE_l+xS_2(q)$,
\begin{equation}\label{eq:id1}
V(x;q)=\sum_lQ_l^2-2x\sum_lg_lQ_l+x^2S_2(q)=\sum_{l=1}^qQ_l^2-x^2S_2(q)-2x\sum_{l=1}^qg_lE_l .
\end{equation}

\emph{Counting pairs.} The sum $\sum_lQ_l^2$ counts ordered pairs $(n,n')$ of squarefree integers in $[1,x]$ with $n\equiv n'\pmod q$. The diagonal contributes $Q(x)$. The pairs with $n<n'$ are the pairs $(n,n+hq)$ with $h\ge1$ and $n+hq\le x$, and the pairs with $n>n'$ are their mirror images. If $hq\ge x$ there is no $n\ge1$ with $n+hq\le x$. Hence
\begin{equation}\label{eq:id2}
\sum_lQ_l^2=Q(x)+2\sum_{1\le h\le x/q}C_x(hq)=Q(x)+2\sum_{1\le h\le x/q}D_x(hq)+2\Sigma_A,
\end{equation}
where $\Sigma_A=\sum_{1\le h\le x/q}A(hq)(x-hq)$.

\emph{The expected pair counts.} Only squarefree $t$ contribute to $A(hq)$, because $\eta$ vanishes elsewhere. For squarefree $t$ and $h\in\N$ we have
\begin{equation}\label{eq:divlemma}
t^2\mid hq\iff m_t(q)\mid h .
\end{equation}
Indeed, with $\gamma=\gcd(t^2,q)$, $t^2\mid hq$ holds if and only if $(t^2/\gamma)\mid h(q/\gamma)$, and since $(t^2/\gamma,q/\gamma)=1$ this holds if and only if $t^2/\gamma=m_t(q)$ divides $h$. Both sums in $\Sigma_A$ are finite ($h\le x/q$ and $t^2\le hq\le x$), so they may be interchanged:
\[
\Sigma_A=A_0\sum_{t\ge1}\eta(t)\sum_{\substack{1\le h\le x/q\\ m_t(q)\mid h}}(x-hq).
\]
Fix $t$ and write $m=m_t(q)$, $M=qm$ and $X=x/M$. Put $h=jm$; then $x-hq=M(X-j)$ and $1\le j\le\lfloor X\rfloor$. With $N=\lfloor X\rfloor$ and $\vartheta=\fr X$,
\[
\sum_{j=1}^{N}(X-j)=NX-\frac{N(N+1)}2=\frac12(X-\vartheta)(X+\vartheta-1)=\frac12\bigl(X^2-X+B(X)\bigr).
\]
This also holds when $N=0$: then $X<1$, $B(X)=X-X^2$, and both sides vanish. Therefore
\begin{equation}\label{eq:id3}
2\Sigma_A=A_0\sum_{t\ge1}\eta(t)\Bigl[\frac{x^2}{q\,m_t(q)}-x+q\,m_t(q)\,B\Bigl(\frac{x}{q\,m_t(q)}\Bigr)\Bigr].
\end{equation}
The bracket vanishes whenever $q\,m_t(q)>x$, since it is twice an empty sum. We now split the bracket into its three terms. Each of the three resulting series converges absolutely: $\eta\ge0$ and $\sum_t\eta(t)<\infty$ by \cref{lem:eta}, $m_t(q)\ge1$, and $0\le q\,m\,B(x/(qm))\le x$. So the splitting is legitimate. By \cref{prop:densities}(iv) and (v), and the definition~\eqref{eq:MT} of $\MT$,
\begin{equation}\label{eq:id4}
2\Sigma_A=x^2S_2(q)-\frac{6x}{\pi^2}+\MT(x;q).
\end{equation}

\emph{Conclusion.} Inserting~\eqref{eq:id4} into~\eqref{eq:id2}, and then~\eqref{eq:id2} into~\eqref{eq:id1}, the terms $x^2S_2(q)$ cancel. This leaves exactly the stated identity. No estimate has been used.
\end{proof}

\begin{remark}
\Cref{prop:identity} says that the size of $V(x;q)$, about $\sqrt{xq}$, comes from the lattice-point discrepancy $B(X)$ in $\sum_{j\le X}(X-j)=\frac12(X^2-X)+\frac12B(X)$. It does not come from any irregularity in the distribution of squarefree numbers. Heuristically, for $t$ coprime to $q$ the term of $\MT(x;q)$ is $A_0q\kappa(t)B\bigl(x/(qt^2)\bigr)$, and the sum of $B\bigl(x/(qt^2)\bigr)$ over $t$ is of order $\sqrt{x/q}$, so $\MT(x;q)$ has order $q\sqrt{x/q}=\sqrt{xq}$; see \cref{sec:sawtooth}. In the computations of \cref{sec:numerics} the identity was also checked numerically: with every other term computed by brute force, it holds at $x=3000$ for all $q\le40$ to a relative accuracy of $\SelfTest$ (\texttt{reproduce.py}, function \texttt{identity\_selftest}).
\end{remark}

\section{Pairs of squarefree numbers with a given difference}\label{sec:pairs}

The following estimate counts pairs of squarefree numbers with a given difference $m$, uniformly in $m\le x$. It is the bottleneck for the lower end of the range in \cref{thm:main}. The exponent $\frac23$ is the one familiar from the closely related additive problem: as reported in~\cite[§1]{BGPVW1998}, an elementary argument of Evelyn and Linfoot gives the number of representations of $n$ as a sum of two $k$-free numbers with an error $O_{\eps,k}(n^{2/(k+1)+\eps})$, which is $O_\eps(n^{2/3+\eps})$ for squarefree numbers. We give a complete proof of the version needed here.

\begin{proposition}[Uniform pair count]\label{prop:pairs}
For every $\eps>0$ there is a constant $c_3(\eps)$ such that
\[
|D_x(m)|\le c_3(\eps)\,x^{2/3+\eps}
\]
for all real $x\ge2$ and all integers $1\le m\le x$, where $D_x(m)$ is defined in~\eqref{eq:pairs-def}.
\end{proposition}

\begin{proof}
Fix $x\ge2$ and $1\le m\le x$, and put $N=x-m\ge0$ and $z=x^{1/3}\ge1$. Then $C_x(m)=\sum_{1\le n\le N}\mu^2(n)\mu^2(n+m)$, with $n$ running over positive integers. Put $\tau^*=\max_{k\le x}\tau(k)$. By \cref{lem:divisor}, $\tau^*\le C(\eps)x^\eps$, and $\tau^*$ also bounds $\tau(m)$.

\emph{Step 1: truncation.} For $k\in\N$ let
\[
S(k)=\sum_{d\le z,\ d^2\mid k}\mu(d),\qquad R(k)=\mu^2(k)-S(k)=\sum_{d>z,\ d^2\mid k}\mu(d).
\]
Then $|S(k)|\le\tau(k)$ and $|R(k)|\le\#\{d>z:\ d^2\mid k\}$. With $\alpha=\mu^2(n)$ and $\beta=\mu^2(n+m)$ we have $\alpha\beta-S(n)S(n+m)=R(n)\beta+S(n)R(n+m)$, and $|\beta|\le1$. Hence
\[
C_x(m)=\sum_{n\le N}S(n)S(n+m)+\mathcal E_1,\qquad|\mathcal E_1|\le\sum_{n\le N}|R(n)|+\tau^*\sum_{n\le N}|R(n+m)|.
\]
For each $d>z$ the number of $n\le N$ with $d^2\mid n$ is at most $N/d^2\le x/d^2$. The number of $n\le N$ with $d^2\mid n+m$ is at most the number of multiples of $d^2$ up to $x$, that is, at most $x/d^2$. By \cref{lem:tail}, both sums are at most $x\sum_{d>z}d^{-2}\le2x/z=2x^{2/3}$. So $|\mathcal E_1|\le2(1+\tau^*)x^{2/3}$.

\emph{Step 2: congruences.} Expanding $S$,
\[
\sum_{n\le N}S(n)S(n+m)=\sum_{d_1,d_2\le z}\mu(d_1)\mu(d_2)\,\#\{1\le n\le N:\ d_1^2\mid n,\ d_2^2\mid n+m\}.
\]
The system $n\equiv0\ (d_1^2)$, $n\equiv-m\ (d_2^2)$ is solvable if and only if $(d_1^2,d_2^2)=(d_1,d_2)^2$ divides $m$. In that case its solutions form one residue class modulo $L=[d_1^2,d_2^2]$, and \cref{lem:class} gives the count $N/L+\theta$. There are at most $\lfloor z\rfloor^2\le x^{2/3}$ pairs $(d_1,d_2)$. So
\[
\sum_{n\le N}S(n)S(n+m)=N\sum_{\substack{d_1,d_2\le z\\(d_1,d_2)^2\mid m}}\frac{\mu(d_1)\mu(d_2)}{[d_1^2,d_2^2]}+\mathcal E_2,\qquad|\mathcal E_2|\le x^{2/3}.
\]

\emph{Step 3: completion.} Let $\Sigma$ be the same sum over all $d_1,d_2\ge1$ with $(d_1,d_2)^2\mid m$. Write $d_1=ea$, $d_2=eb$ with $e=(d_1,d_2)$ and $(a,b)=1$; then $[d_1^2,d_2^2]=e^2a^2b^2$, and the condition reads $e^2\mid m$. By symmetry, the pairs with $\max(d_1,d_2)>z$ contribute in absolute value at most twice the contribution of the pairs with $d_1>z$, that is, at most
\[
2\sum_{e^2\mid m}\frac1{e^2}\sum_{a>z/e}\frac1{a^2}\sum_{b\ge1}\frac1{b^2}\le2\cdot\frac{\pi^2}6\sum_{e^2\mid m}\frac1{e^2}\cdot\frac{2e}{z}=\frac{2\pi^2}{3z}\sum_{e^2\mid m}\frac1e\le\frac{7\,\tau(m)}{z},
\]
by \cref{lem:tail} with $w=z/e$. This also shows that $\Sigma$ converges absolutely. Hence
\[
N\sum_{\substack{d_1,d_2\le z\\(d_1,d_2)^2\mid m}}\frac{\mu(d_1)\mu(d_2)}{[d_1^2,d_2^2]}=N\Sigma+\mathcal E_3,\qquad|\mathcal E_3|\le\frac{7x\tau^*}{z}=7\tau^*x^{2/3}.
\]

\emph{Step 4: evaluation of $\Sigma$.} A nonzero term needs $d_1=ea$ and $d_2=eb$ squarefree. Then $e$, $a$, $b$ are squarefree and pairwise coprime, and $\mu(ea)\mu(eb)=\mu(a)\mu(b)$. By absolute convergence and \cref{lem:pairs-euler},
\[
\Sigma=\sum_{\substack{e\ \mathrm{squarefree}\\ e^2\mid m}}\frac1{e^2}\sum_{\substack{\mu^2(ab)=1\\(ab,e)=1}}\frac{\mu(a)\mu(b)}{a^2b^2}
=A_0\sum_{\substack{e\ \mathrm{squarefree}\\ e^2\mid m}}\frac1{e^2}\prod_{p\mid e}\frac{p^2}{p^2-2}=A_0\sum_{t^2\mid m}\eta(t)=A(m).
\]

\emph{Conclusion.} Since $N=x-m$, we get $C_x(m)=A(m)(x-m)+\mathcal E_1+\mathcal E_2+\mathcal E_3$, so
\[
|D_x(m)|\le\bigl(2(1+\tau^*)+1+7\tau^*\bigr)x^{2/3}\le(3+9C(\eps))\,x^{2/3+\eps}.\qedhere
\]
\end{proof}

\section{Sawtooth sums and the main term for one modulus}\label{sec:sawtooth}

\subsection{Sawtooth sums}

The function $B(u)=\fr u(1-\fr u)$ is continuous, $1$-periodic, takes values in $[0,\frac14]$, and satisfies $B(u)=u-u^2$ for $0\le u\le1$. The sums we need are $\sum_{s\ge1}B(Z/s^2)$ and weighted versions of them. Their size $\asymp\sqrt Z$ has two sources: the terms with $s>\sqrt Z$, where $B(Z/s^2)=Z/s^2-Z^2/s^4$, and the terms with $s\le\sqrt Z$, where $B$ oscillates. The natural approximation is the integral $\int_0^\infty B(Z/u^2)\,du=I\sqrt Z$ (substitute $u=\sqrt Z\,w$). For $q\in\N$ put
\[
\kappa_q(t)=\kappa(t)\one[(t,q)=1],\qquad \Delta_q=\prod_p\Bigl(1-\frac1p\Bigr)\Bigl(1+\one[p\nmid q]\frac{p}{p^2-2}\Bigr).
\]
For $p\nmid q$ the factor of $\Delta_q$ is $(p-1)^2(p+2)/(p(p^2-2))=1+O(p^{-2})$; there are finitely many other factors, so the product converges absolutely.

\begin{proposition}[Sawtooth sums]\label{prop:sawtooth}
\leavevmode
\begin{enumerate}[label=\textup{(\alph*)}]
\item $I=\int_0^\infty B(u^{-2})\,du=\frac12\int_0^\infty B(v)v^{-3/2}\,dv=-4\zeta(-\tfrac12)$.
\item For every real $Z>0$, $\bigl|\sum_{s\ge1}B(Z/s^2)-I\sqrt Z\bigr|\le c_1Z^{1/3}$ with $c_1=60$.
\item For every $\eps>0$ there is $c_5(\eps)$ such that for all $q\in\N$ and real $Z>0$,
\[
\Bigl|\sum_{t\ge1}\kappa_q(t)B(Z/t^2)-\Delta_q\,I\sqrt Z\Bigr|\le c_5(\eps)\,q^\eps Z^{1/3}.
\]
\end{enumerate}
\end{proposition}

\begin{proof}
(a) The substitution $v=u^{-2}$ gives the second expression. The integral converges, because $B(v)\le v$ near $0$ and $B(v)\le\frac14$. Since $B$ is continuous and piecewise $C^1$ with $B'(v)=1-2\fr v$, integration by parts on $[0,X]$ gives
\[
\int_0^XB(v)v^{-3/2}\,dv=\bigl[-2B(v)v^{-1/2}\bigr]_0^X+2\int_0^XB'(v)v^{-1/2}\,dv .
\]
The boundary term tends to $0$, because $B(v)v^{-1/2}\le v^{1/2}$ near $0$ and $B(X)X^{-1/2}\le\frac14X^{-1/2}$. Hence $\int_0^\infty Bv^{-3/2}\,dv=2\lim_{N\to\infty}\int_0^NB'v^{-1/2}\,dv$ along integers $N$. On $(n,n+1)$ we have $B'(n^+)=1$, $B'((n+1)^-)=-1$ and $B''=-2$, so integrating by parts with the antiderivative $2v^{1/2}$ of $v^{-1/2}$,
\[
\int_n^{n+1}B'(v)v^{-1/2}\,dv=-2(n+1)^{1/2}-2n^{1/2}+4\int_n^{n+1}v^{1/2}\,dv,
\]
and summing over $0\le n<N$,
\begin{equation}\label{eq:Bprime}
\int_0^NB'(v)v^{-1/2}\,dv=-4\sum_{m=1}^Nm^{1/2}+2N^{1/2}+\tfrac83N^{3/2}.
\end{equation}
Now apply \cref{lem:EM} on $(N,\infty)$ to $G(u)=u^{-s}$ with $\Re s>1$, for which $\sum_{n>N}n^{-s}$ and $\int_N^\infty u^{-s}\,du=N^{1-s}/(s-1)$ converge absolutely, using $\psi(N)=-\frac12$ and $\psi_2(N)=\frac1{12}$:
\[
\sum_{n>N}n^{-s}=\frac{N^{1-s}}{s-1}-\frac{N^{-s}}2+\frac s{12}N^{-s-1}-s(s+1)\int_N^\infty\psi_2(u)u^{-s-2}\,du .
\]
The right-hand side is analytic in $\Re s>-1$, $s\ne1$, because the integral converges locally uniformly there. By uniqueness of analytic continuation, $\zeta(s)=\sum_{n\le N}n^{-s}$ plus the right-hand side in that region. At $s=-\frac12$ we have $s(s+1)=-\frac14$ and $|\int_N^\infty\psi_2u^{-3/2}\,du|\le\frac16N^{-1/2}$, so
\[
\sum_{m\le N}m^{1/2}=\zeta(-\tfrac12)+\tfrac23N^{3/2}+\tfrac12N^{1/2}+O(N^{-1/2}).
\]
Inserting this into~\eqref{eq:Bprime} shows that $\int_0^NB'v^{-1/2}\,dv\to-4\zeta(-\frac12)$. Hence $\int_0^\infty Bv^{-3/2}\,dv=-8\zeta(-\frac12)$ and $I=-4\zeta(-\frac12)$. By the functional equation~\cite[Corollary~10.4]{MV2007} at $s=-\frac12$, $\zeta(-\frac12)=2^{-1/2}\pi^{-3/2}\Gamma(\frac32)\sin(-\frac\pi4)\zeta(\frac32)=-\zeta(\frac32)/(4\pi)$. So $I=\zeta(3/2)/\pi>0$, the value stated in~\eqref{eq:I}.

(b) \emph{The case $Z<1$.} Here $0\le B(Z/s^2)\le Z/s^2$, so $0\le\sum_sB(Z/s^2)\le\zeta(2)Z<2Z^{1/3}$. Also $0<I\sqrt Z<Z^{1/3}$, since $I<1$. The difference of two numbers in $[0,2Z^{1/3}]$ is at most $2Z^{1/3}$.

\emph{The case $Z\ge1$.} Put $\Lambda=Z^{1/3}\ge1$ and $T=\lfloor\Lambda\rfloor\ge1$. Then $\Lambda/2<T\le\Lambda$, hence
\begin{equation}\label{eq:b1}
\frac Z{T^3}<8,\qquad K_0:=\frac Z{T^2}<4\Lambda,\qquad \frac{Z^2}{T^5}<32\Lambda .
\end{equation}
\emph{Small $s$.} Both $\sum_{s\le T}B(Z/s^2)$ and $\int_0^TB(Z/u^2)\,du$ lie in $[0,T/4]\subset[0,\Lambda/4]$, so their difference is at most $\Lambda/4$. It therefore suffices to show
\begin{equation}\label{eq:b2}
\Bigl|\sum_{s>T}B(Z/s^2)-\int_T^\infty B(Z/u^2)\,du\Bigr|\le57\Lambda ,
\end{equation}
which gives the claim with $57+\frac14<60$.

\emph{A decomposition of $B$.} For real $w\ge0$ with $j=\lfloor w\rfloor$ we have $B(w)=(w-j)(1-w+j)=w-w^2+2jw-j-j^2$, and $\sum_{k=1}^j2(w-k)=2jw-j^2-j$. Hence
\[
B(w)=w-w^2+2\sum_{k\ge1}(w-k)_+,\qquad(v)_+=\max(v,0).
\]
Put $\Phi(u)=Z/u^2-Z^2/u^4$ and $G_k(u)=(Z/u^2-k)_+$, so that $B(Z/u^2)=\Phi(u)+2\sum_{k\ge1}G_k(u)$. For $u\ge T$ only the $k<K_0$ contribute, and there are fewer than $K_0<4\Lambda$ of them. For such $k$, $G_k(u)=Z/u^2-k$ on $[T,R_k]$ with $R_k=\sqrt{Z/k}>T$, and $G_k=0$ on $[R_k,\infty)$.

\emph{Euler--Maclaurin.} Apply \cref{lem:EM} to $\Phi$ on $(T,\infty)$ (where $\sum_{s>T}\Phi(s)$ and $\int_T^\infty\Phi$ converge absolutely, since $\Phi(u)=O(u^{-2})$), and to $Z/u^2-k$ on $(T,R_k]$ for each $k<K_0$. Use $\psi(T)=-\frac12$, $\psi_2(T)=\frac1{12}$ and $G_k(R_k)=0$, and add. Since $\int_T^\infty\Phi+2\sum_k\int_T^{R_k}G_k=\int_T^\infty B(Z/u^2)\,du$ (all integrals converge absolutely), this gives
\[
\sum_{s>T}B(Z/s^2)=\int_T^\infty B(Z/u^2)\,du+R_1+R_2+R_3+R_4+R_5,
\]
where, writing $G_k'$ and $G_k''$ for the derivatives of $Z/u^2-k$:
\begin{itemize}
\item $R_1=\psi(T)\bigl[\Phi(T)+2\sum_{k<K_0}G_k(T)\bigr]=\psi(T)B(Z/T^2)$. The individual terms here have size about $Z^2/T^4\asymp Z^{2/3}$, but their sum is $B(Z/T^2)\le\frac14$. So $|R_1|\le\frac18$. This is the one place where cancellation is essential.
\item $R_2=-\frac1{12}\bigl[\Phi'(T)+2\sum_{k<K_0}G_k'(T)\bigr]$, with $\Phi'(T)=-2Z/T^3+4Z^2/T^5$ and $G_k'(T)=-2Z/T^3$. By~\eqref{eq:b1}, $|\Phi'(T)|<16+128\Lambda$ and $|2\sum_kG_k'(T)|\le4K_0Z/T^3<128\Lambda$. So $|R_2|\le\frac1{12}(16+256\Lambda)\le\frac{272}{12}\Lambda$.
\item $R_3=-\int_T^\infty\psi_2\Phi''\,du$ with $\Phi''=6Z/u^4-20Z^2/u^6$. So $|R_3|\le\frac1{12}(2Z/T^3+4Z^2/T^5)<\frac1{12}(16+128\Lambda)\le12\Lambda$.
\item $R_4=-2\sum_{k<K_0}\int_T^{R_k}\psi_2G_k''\,du$ with $G_k''=6Z/u^4$ and $\int_T^\infty6Zu^{-4}\,du=2Z/T^3<16$. So $|R_4|\le\frac2{12}K_0\cdot16<\frac{32}3\Lambda$.
\item $R_5=2\sum_{k<K_0}\psi_2(R_k)G_k'(R_k)$ with $G_k'(R_k)=-2Z/R_k^3=-2k^{3/2}/\sqrt Z$. Since $k<4\Lambda$, $\sum_{k<K_0}k^{3/2}<4\Lambda\cdot8\Lambda^{3/2}$, and with $\sqrt Z=\Lambda^{3/2}$ we get $|R_5|\le\frac2{12}\cdot2\cdot32\Lambda=\frac{32}3\Lambda$. These are the kinks of the $G_k$ at the points $R_k$. Summing over the $K_0\asymp Z^{1/3}$ values of $k$ is what produces the exponent $\frac13$.
\end{itemize}
Using $\Lambda\ge1$, the total is at most $(\frac18+\frac{272}{12}+12+\frac{32}3+\frac{32}3)\Lambda=\frac{449}8\Lambda<57\Lambda$, which proves~\eqref{eq:b2}.

(c) Fix $q$. The function $\kappa_q$ is multiplicative with $\kappa_q(p)=p^2/(p^2-2)$ for $p\nmid q$, $\kappa_q(p)=0$ for $p\mid q$ and $\kappa_q(p^j)=0$ for $j\ge2$. Let $\beta=\mu*\kappa_q$, so that $\kappa_q=1*\beta$. Then $\beta$ is multiplicative, with $\beta(p^j)=\kappa_q(p^j)-\kappa_q(p^{j-1})$:
\begin{itemize}
\item for $p\mid q$: $\beta(p)=-1$ and $\beta(p^j)=0$ for $j\ge2$;
\item for $p\nmid q$: $\beta(p)=\frac2{p^2-2}$, $\beta(p^2)=-\frac{p^2}{p^2-2}$ and $\beta(p^j)=0$ for $j\ge3$.
\end{itemize}
\emph{Absolute convergence.} For $p\nmid q$, $\sum_j|\beta(p^j)|p^{-2j/3}\le\frac{2p^{-2/3}}{p^2-2}+2p^{-4/3}$, because $p^2/(p^2-2)\le2$. Hence
\[
\Sigma_{2/3}:=\sum_{n\ge1}\frac{|\beta(n)|}{n^{2/3}}\le\prod_{p\mid q}\bigl(1+p^{-2/3}\bigr)\prod_p\Bigl(1+\frac{2p^{-2/3}}{p^2-2}+2p^{-4/3}\Bigr)\le C_\beta\,2^{\omega(q)}\le C_\beta C(\eps)q^\eps,
\]
with an absolute constant $C_\beta$, by \cref{lem:divisor}. Similarly $\sum_n|\beta(n)|/n<\infty$ for each fixed $q$.

\emph{Rearrangement.} For $W>0$, $\sum_mB(W/m^2)\le\sum_m\min(\frac14,W/m^2)\le3\sqrt W$. Indeed, for $W<1$ the sum is at most $\zeta(2)W\le2\sqrt W$, and for $W\ge1$ it is at most $\frac14\sqrt W+W\cdot2/\sqrt W$ by \cref{lem:tail}. Hence $\sum_{n,m}|\beta(n)|B(Z/(nm)^2)\le3\sqrt Z\sum_n|\beta(n)|/n<\infty$, and we may rearrange:
\[
\sum_{t\ge1}\kappa_q(t)B(Z/t^2)=\sum_{n,m\ge1}\beta(n)B\Bigl(\frac{Z/n^2}{m^2}\Bigr)=\sum_{n\ge1}\beta(n)\sum_{m\ge1}B\Bigl(\frac{Z/n^2}{m^2}\Bigr).
\]
\emph{Main term.} Part (b) holds for every positive argument. Applying it with $Z/n^2$ in place of $Z$, and using $I\sqrt{Z/n^2}=I\sqrt Z/n$,
\[
\Bigl|\sum_t\kappa_q(t)B(Z/t^2)-I\sqrt Z\sum_n\frac{\beta(n)}n\Bigr|\le c_1\sum_n|\beta(n)|\Bigl(\frac Z{n^2}\Bigr)^{1/3}=c_1\Sigma_{2/3}Z^{1/3}.
\]
\emph{The constant.} By~\eqref{eq:euler}, $\sum_n\beta(n)/n=\prod_p\bigl(1+\beta(p)/p+\beta(p^2)/p^2\bigr)$. For $p\mid q$ the factor is $1-1/p$. For $p\nmid q$ it is
\[
1+\frac2{p(p^2-2)}-\frac1{p^2-2}=\frac{p^3-3p+2}{p(p^2-2)}=\frac{(p-1)^2(p+2)}{p(p^2-2)}=\Bigl(1-\frac1p\Bigr)\Bigl(1+\frac p{p^2-2}\Bigr).
\]
Hence $\sum_n\beta(n)/n=\Delta_q$, and (c) holds with $c_5(\eps)=c_1C_\beta C(\eps)$.
\end{proof}

\subsection{The main term for one modulus}

\begin{proposition}[Asymptotics of $\MT$]\label{prop:MT}
For every $\eps>0$ there is $c_6(\eps)$ such that for all real $x>0$ and all $q\in\N$,
\[
\bigl|\MT(x;q)-h(q)\sqrt{xq}\bigr|\le c_6(\eps)\,x^{1/3}q^{2/3+\eps},
\]
with $h$ as in~\eqref{eq:h}. In particular this holds for $x\ge1$ and $1\le q\le x$.
\end{proposition}

\begin{proof}
\emph{Factorising $t$.} Fix $q$. Every squarefree $t$ factors uniquely as $t=t_0t_1t_2$, where $t_0$ is composed of primes $p\nmid q$, $t_1$ of primes $p\,\|\,q$, and $t_2$ of primes $p$ with $p^2\mid q$. Conversely, every such triple of squarefree numbers gives a squarefree $t$. Let $\mathcal S_1$ and $\mathcal S_2$ be the finite sets of admissible $t_1$ and $t_2$ (all of them divide $q$). For squarefree $t$, $\gcd(t^2,q)=\prod_{p\mid t}p^{\min(2,v_p(q))}=t_1t_2^2$, so
\[
m_t(q)=t_0^2t_1,\qquad\frac x{q\,m_t(q)}=\frac{Z_1}{t_0^2}\quad\text{with}\quad Z_1=\frac{x}{q\,t_1},\qquad \eta(t)m_t(q)=\eta(t_1)\eta(t_2)\,t_1\,\kappa(t_0).
\]
By the absolute convergence in \cref{lem:eta} we may regroup:
\[
\MT(x;q)=A_0\,q\sum_{t_1\in\mathcal S_1}\sum_{t_2\in\mathcal S_2}\eta(t_1)\eta(t_2)\,t_1\,\Sigma(t_1),\qquad\Sigma(t_1)=\sum_{t_0\ge1}\kappa_q(t_0)B\Bigl(\frac{Z_1}{t_0^2}\Bigr).
\]
\emph{Applying \cref{prop:sawtooth}(c).} With $Z=Z_1>0$ we get $\Sigma(t_1)=\Delta_qI\sqrt{Z_1}+\theta\,c_5(\eps)q^\eps Z_1^{1/3}$. Since $t_1\sqrt{Z_1}=\sqrt{x/q}\,t_1^{1/2}$ and $t_1Z_1^{1/3}=(x/q)^{1/3}t_1^{2/3}$,
\[
\MT(x;q)=A_0\Delta_qI\sqrt{xq}\sum_{t_1\in\mathcal S_1}\eta(t_1)t_1^{1/2}\sum_{t_2\in\mathcal S_2}\eta(t_2)+R,
\]
\[
|R|\le c_5(\eps)\,q^{1+\eps}\Bigl(\frac xq\Bigr)^{1/3}\sum_{t_1\in\mathcal S_1}\eta(t_1)t_1^{2/3}\sum_{t_2\in\mathcal S_2}\eta(t_2),
\]
where we used $A_0\le1$. By~\eqref{eq:eta-bound}, $\sum_{t_1}\eta(t_1)t_1^{2/3}\le\prod_p(1+2p^{-4/3})$ and $\sum_{t_2}\eta(t_2)\le\prod_p(1+2p^{-2})$, both absolute constants. So $|R|\le c_6(\eps)x^{1/3}q^{2/3+\eps}$.

\emph{The constant is $h(q)$.} The two finite sums are Euler products over the primes dividing $q$:
\[
\sum_{t_1\in\mathcal S_1}\eta(t_1)t_1^{1/2}=\prod_{p\|q}\Bigl(1+\frac{\sqrt p}{p^2-2}\Bigr),\qquad\sum_{t_2\in\mathcal S_2}\eta(t_2)=\prod_{p^2\mid q}\Bigl(1+\frac1{p^2-2}\Bigr).
\]
Multiplying the absolutely convergent products $A_0$ and $\Delta_q$ factor by factor, the local factor at $p$ of the main coefficient is
\begin{align*}
p\nmid q:&\quad\Bigl(1-\tfrac2{p^2}\Bigr)\Bigl(1-\tfrac1p\Bigr)\Bigl(1+\tfrac p{p^2-2}\Bigr)=\frac{(p-1)(p^2+p-2)}{p^3}=\frac{(p-1)^2(p+2)}{p^3}=a_p,\\
p\,\|\,q:&\quad\Bigl(1-\tfrac2{p^2}\Bigr)\Bigl(1-\tfrac1p\Bigr)\Bigl(1+\tfrac{\sqrt p}{p^2-2}\Bigr)=\frac{(p-1)(p^2-2+\sqrt p)}{p^3}=b_p,\\
p^2\mid q:&\quad\Bigl(1-\tfrac2{p^2}\Bigr)\Bigl(1-\tfrac1p\Bigr)\Bigl(1+\tfrac1{p^2-2}\Bigr)=\frac{(p-1)(p^2-1)}{p^3}=c_p .
\end{align*}
Together with the factor $I$, this is exactly $h(q)$.
\end{proof}

\begin{remark}
Since $b_p/a_p\ge1-1/p$ and $c_p/a_p=(p+1)/(p+2)\ge1-1/p$, we have $h(q)\ge IA_h\,\varphi(q)/q\gg_\eps q^{-\eps}$, where $A_h=\prod_pa_p>0$ (\cref{prop:poly}(0)). So \cref{prop:MT} gives, for $x\ge1$ and $1\le q\le x$,
\[
\MT(x;q)=h(q)\sqrt{xq}\,\bigl(1+O_\eps(x^{2\eps}(q/x)^{1/6})\bigr),
\]
because the relative error is $\ll_\eps q^{2\eps}(q/x)^{1/6}$ and $q^{2\eps}\le x^{2\eps}$ for $q\le x$. (Without the restriction $q\le x$ the display can fail. Since $0\le q\,m\,B(x/(qm))\le x$ and $\sum_t\eta(t)<\infty$ (\cref{lem:eta}), \eqref{eq:MT} gives $0\le\MT(x;q)\ll x$ for all $x>0$. Hence $\MT(x;1)/(h(1)\sqrt x)\to0$ as $x\to0$, whereas for $\eps>1/12$ the displayed error term $x^{2\eps}x^{-1/6}$ tends to $0$.) The relative error is not small when $q\asymp x$. This is why the asymptotic formula of \cref{thm:main}(iii) requires $Y\le x^{1-\delta}$, whereas the error term in~(ii), which keeps $\MT$ unexpanded, has no such restriction. The numerical data in \cref{sec:numerics} suggest that, after summation over the moduli, the secondary term has relative size roughly proportional to $\sqrt{Y/x}$. That is consistent with, but sharper than, the proved bound.
\end{remark}

\section{The variance for a single modulus}\label{sec:single}

\begin{proposition}[Single-modulus variance]\label{prop:single}
For every $\eps>0$ there is $c_7(\eps)$ such that for all real $x\ge2$ and all integers $1\le q\le x$:
\begin{enumerate}[label=\textup{(\roman*)}]
\item $|V(x;q)-\MT(x;q)|\le c_7(\eps)\bigl(x^{5/3+\eps}/q+\sqrt x\bigr)$;
\item $|V(x;q)-h(q)\sqrt{xq}|\le c_7(\eps)\bigl(x^{1/3}q^{2/3+\eps}+x^{5/3+\eps}/q+\sqrt x\bigr)$.
\end{enumerate}
\end{proposition}

\begin{proof}
(i) By \cref{prop:identity}, which holds for all $x\ge1$ and $q\in\N$,
\[
V(x;q)-\MT(x;q)=2\sum_{1\le h\le x/q}D_x(hq)+\Bigl(Q(x)-\frac{6x}{\pi^2}\Bigr)-2x\sum_{l=1}^qg(q,l)E(x;q,l).
\]
For $1\le h\le x/q$ the integer $m=hq$ satisfies $1\le m\le x$. Since $x\ge2$, \cref{prop:pairs} applies to each of the at most $x/q$ terms. (This is the only place where $x\ge2$ rather than $x\ge1$ is needed.) It gives
\[
\Bigl|2\sum_{1\le h\le x/q}D_x(hq)\Bigr|\le2\cdot\frac xq\cdot c_3(\eps)x^{2/3+\eps}=2c_3(\eps)\frac{x^{5/3+\eps}}q .
\]
\Cref{prop:trivial}(i) gives $|Q(x)-6x/\pi^2|\le c_0\sqrt x$. \Cref{prop:trivial}(ii) gives
\[
\Bigl|2x\sum_lg(q,l)E(x;q,l)\Bigr|\le2c_0\frac{x^{3/2}}q\le2c_0\frac{x^{5/3+\eps}}q,
\]
since $x\ge1$. By the triangle inequality, (i) holds with $c_7(\eps)\ge2c_3(\eps)+2c_0$.

(ii) \Cref{prop:MT} gives $|\MT(x;q)-h(q)\sqrt{xq}|\le c_6(\eps)x^{1/3}q^{2/3+\eps}$. Combining this with (i) by the triangle inequality proves (ii) with
\[
c_7(\eps)=\max\bigl\{2c_3(\eps)+2c_0,\ c_6(\eps)\bigr\},
\]
and (i) remains true with this larger constant.
\end{proof}

\section{Averages over polynomial moduli}\label{sec:poly}

In this section $f$ is fixed and satisfies~\eqref{eq:f-hyp}. All constants may depend on $f$, and on $\eps$ where indicated. Write $f(t)=\alpha t^d+\dots$, with leading coefficient $\alpha\ge1$, and let $P_0$ be the finite set of primes dividing the content of $f$ (the gcd of its coefficients). Put $x_p=\rho(p)/p$ and $z_p=\rho(p^2)/p^2$.

\begin{proposition}[Polynomial averages]\label{prop:poly}
\leavevmode
\begin{enumerate}[label=\textup{(\arabic*)},start=0]
\item $\sigma_p>0$ for every $p$, and $\prod_p\sigma_p$ converges absolutely. Hence $0<C_f<\infty$.
\end{enumerate}
Moreover, for every $\eps>0$ there are constants $c_8(f)$ and $c_8(f,\eps)$ such that for all real $y\ge1$, with $Y=f(y)$:
\begin{enumerate}[label=\textup{(\alph*)}]
\item $\displaystyle\Bigl|\sum_{1\le k\le y}F(k)\,h(f(k))-C_fY^{3/2}\Bigr|\le c_8(f,\eps)\,Y^{3/2+\eps}y^{-1/2}$;
\item $\displaystyle\sum_{1\le k\le y}\frac{f'(k)}{f(k)}\le c_8(f)\,(1+\log Y)$;
\item $\displaystyle\sum_{1\le k\le y}f'(k)\le c_8(f)\,Y$.
\end{enumerate}
\end{proposition}

The main difficulty in~(a) is that $f$ may have repeated factors or a nontrivial content. Then $\rho(p^2)$ may be as large as $p\rho(p)$, and $\rho(p)=p$ is possible for $p\in P_0$. The proof uses only the bounds of \cref{lem:rho} and never assumes $\rho(p^2)=\rho(p)$.

\subsection{Preliminaries}

\begin{lemma}[Local factors]\label{lem:local}
For every prime $p$, $a_p,b_p,c_p>0$, and
\[
b_p-a_p=-\frac1p+p^{-3/2}+p^{-2}-p^{-5/2},\qquad c_p-b_p=-p^{-3/2}\bigl(1-p^{-1/2}\bigr)\bigl(1-p^{-1}\bigr).
\]
Consequently $a_p\ge\frac12$, $0<1-a_p\le3/p^2$, $|b_p-a_p|<1/p$ and $-p^{-3/2}<c_p-b_p<0$.
\end{lemma}

\begin{proof}
Positivity is clear from~\eqref{eq:abc}. Expanding~\eqref{eq:abc} gives $b_p=1-\frac1p-\frac2{p^2}+\frac2{p^3}+p^{-3/2}-p^{-5/2}$ and $c_p=1-\frac1p-\frac1{p^2}+\frac1{p^3}$, from which the two identities follow by subtraction. (They were also confirmed with computer algebra; see \cref{sec:numerics}.) Next, $a_2=\frac12$, and $a_p\ge1-3/p^2\ge\frac23$ for $p\ge3$. The difference $1-a_p=3/p^2-2/p^3$ lies in $(0,3/p^2]$. In the first identity $-1/p+p^{-3/2}\in(-1/p,0)$ and $p^{-2}-p^{-5/2}\in[0,p^{-2})$, so $|b_p-a_p|<1/p$. In the second, $(1-p^{-1/2})(1-p^{-1})\in(0,1)$.
\end{proof}

\begin{lemma}[Root counts]\label{lem:rho}
The function $\rho$ is multiplicative. For every prime $p$, $\rho(p^2)\le p\,\rho(p)$ and $\rho(p^j)\le p^j$. For $p\notin P_0$, $\rho(p)\le d$. Hence $0\le z_p\le x_p\le1$ for every $p$, and $x_p,z_p\le d/p$ for $p\notin P_0$.
\end{lemma}

\begin{proof}
Multiplicativity follows from the Chinese remainder theorem. Every root of $f$ modulo $p^2$ reduces to a root modulo $p$, and each residue modulo $p$ has exactly $p$ lifts modulo $p^2$, so $\rho(p^2)\le p\rho(p)$. For $p\notin P_0$ the reduction of $f$ modulo $p$ is a nonzero polynomial of degree at most $d$ over the field $\Z/p\Z$, so it has at most $d$ roots.
\end{proof}

\begin{lemma}[Shape of $F$]\label{lem:shape}
There are constants $0<\gamma_1\le1\le\gamma_2$, depending only on $f$, such that for all real $t\ge1$
\begin{equation}\label{eq:shape}
\gamma_1t^d\le f(t)\le\gamma_2t^d,\qquad\gamma_1t^{d-1}\le f'(t)\le\gamma_2t^{d-1},\qquad\gamma_1t^{3d/2-1}\le F(t)\le\gamma_2t^{3d/2-1}.
\end{equation}
Consequently, for $y\ge1$ and $Y=f(y)$: $\sup_{[1,y]}F\le\gamma_3F(y)$ with $\gamma_3=\gamma_2/\gamma_1$; $F(y)\le\gamma_4Y^{3/2}/y$ with $\gamma_4=\gamma_2\gamma_1^{-3/2}$; $y\le Y/\gamma_1$; and $\int_1^yF(t)\,dt=\frac23\bigl(Y^{3/2}-f(1)^{3/2}\bigr)$. Finally, $[1,\infty)$ is the union of at most $2d-1$ closed intervals, on each of which $F$ is monotone.
\end{lemma}

\begin{proof}
On $[1,\infty)$ the functions $f(t)/t^d$, $f'(t)/t^{d-1}$ and $F(t)/t^{3d/2-1}$ are continuous and positive, by~\eqref{eq:f-hyp}. They have the positive limits $\alpha$, $\alpha d$ and $\alpha^{3/2}d$ as $t\to\infty$, so they are bounded above and below by positive constants. Since $3d/2-1>0$, $\sup_{[1,y]}F\le\gamma_2y^{3d/2-1}\le\gamma_3F(y)$. Also $Y^{3/2}/y\ge\gamma_1^{3/2}y^{3d/2-1}\ge\gamma_1^{3/2}\gamma_2^{-1}F(y)$, and $Y\ge\gamma_1y^d\ge\gamma_1y$. The integral follows from $(f^{3/2})'=\frac32f'f^{1/2}$. Finally $F'=P/(2\sqrt f)$ with $P=2ff''+f'^2$. This is a polynomial of degree $2d-2$ with leading coefficient $\alpha^2d(3d-2)>0$ (for $d=1$, $P=\alpha^2$), so it has at most $2d-2$ real zeros, and between consecutive zeros $F'$ has constant sign.
\end{proof}

\begin{lemma}[Monotone weights in a residue class]\label{lem:monotone}
Let $J\subset\R$ be a bounded interval (open, closed or half-open) and $G\ge0$ monotone on $J$. Let $e\in\N$ and $r\in\Z$, and let $k_1<\dots<k_n$ be the integers $k\in J$ with $k\equiv r\pmod e$. Then
\[
\Bigl|\sum_{j=1}^nG(k_j)-\frac1e\int_JG(t)\,dt\Bigr|\le3\sup_JG .
\]
\end{lemma}

\begin{proof}
If $n=0$, then $J$ has length at most $e$ and $\frac1e\int_JG\le\sup_JG$. Let $n\ge1$ and suppose $G$ is nonincreasing. Then $G(k_{j+1})\le\frac1e\int_{k_j}^{k_{j+1}}G\le G(k_j)$ for $j<n$. Summing over $j$, $\sum_{j\ge2}G(k_j)\le\frac1e\int_{k_1}^{k_n}G\le\sum_{j<n}G(k_j)$, so $\bigl|\sum_jG(k_j)-\frac1e\int_{k_1}^{k_n}G\bigr|\le\sup_JG$. Since $k_1-e\notin J$ and $k_n+e\notin J$, the parts of $J$ to the left of $k_1$ and to the right of $k_n$ have length at most $e$. Each contributes at most $\sup_JG$ to $\frac1e\int_JG$. The nondecreasing case is symmetric.
\end{proof}

The lemma has content only when $\sup_JG<\infty$. In its application below, $G=F$ is continuous on the compact interval $[1,y]$, so this holds.

\subsection{Proof of \texorpdfstring{\cref{prop:poly}}{Proposition 8.1}(0), (b) and (c)}

(0) By \cref{lem:rho} the weights $1-x_p$, $x_p-z_p$ and $z_p$ in~\eqref{eq:sigma} are nonnegative and sum to $1$. So $\sigma_p$ is a convex combination of $a_p,b_p,c_p>0$, and $\sigma_p\ge\min(a_p,b_p,c_p)>0$. Also $\sigma_p=a_p+(b_p-a_p)x_p+(c_p-b_p)z_p$. For $p\notin P_0$, \cref{lem:local,lem:rho} give
\begin{equation}\label{eq:sigma-close}
|\sigma_p-1|\le(1-a_p)+\frac1p\cdot\frac dp+p^{-3/2}\cdot\frac dp\le\frac{3+2d}{p^2}.
\end{equation}
Hence $\sum_p|\sigma_p-1|<\infty$. Since every $\sigma_p>0$, the product $\prod_p\sigma_p$ converges absolutely to a positive finite number. As $I>0$ (\cref{prop:sawtooth}(a)), $0<C_f<\infty$. In the same way $A_h:=\prod_pa_p$ converges absolutely, and $0<A_h<1$.

(b), (c) By~\eqref{eq:shape}, for $1\le k\le y$ we have $f'(k)/f(k)\le\gamma_3/k$ and $f'(k)\le\gamma_2k^{d-1}$. Hence
\[
\sum_{k\le y}\frac{f'(k)}{f(k)}\le\gamma_3(1+\log y)\le\gamma_3\bigl(1+\log Y+\log(1/\gamma_1)\bigr)\le c_8(f)(1+\log Y),
\]
using $y\le Y/\gamma_1$. Also $\sum_{k\le y}f'(k)\le\gamma_2y\cdot y^{d-1}\le\gamma_3Y$.

\subsection{Proof of \texorpdfstring{\cref{prop:poly}}{Proposition 8.1}(a)}

\emph{A divisor-sum expansion of $h$.} By~\eqref{eq:h}, and since $\prod_pa_p$ converges absolutely,
\[
h(q)=I\,A_h\,\hnat(q),\qquad\hnat(q)=\prod_{p\|q}\frac{b_p}{a_p}\prod_{p^2\mid q}\frac{c_p}{a_p}.
\]
The function $\hnat$ is multiplicative, with $\hnat(p)=b_p/a_p$ and $\hnat(p^j)=c_p/a_p$ for $j\ge2$. Let $w=\mu*\hnat$, so that $\hnat(q)=\sum_{e\mid q}w(e)$. Then $w$ is multiplicative with
\[
w(p)=\frac{b_p-a_p}{a_p},\qquad w(p^2)=\frac{c_p-b_p}{a_p},\qquad w(p^j)=0\quad(j\ge3).
\]
By \cref{lem:local} and $a_p\ge\frac12$,
\begin{equation}\label{eq:w-local}
|w(p)|\le\frac2p,\qquad|w(p^2)|\le2p^{-3/2}.
\end{equation}
Thus $w(e)=0$ unless $e$ is cube-free, and every cube-free $e$ is uniquely $e=ab^2$ with $a,b$ squarefree and coprime. Since $2/p\le4p^{-3/4}$ and $2p^{-3/2}=2(p^2)^{-3/4}$,
\begin{equation}\label{eq:w-bounds}
|w(e)|\le4^{\omega(e)}e^{-3/4},\qquad |w(ab^2)|\le2^{\omega(ab)}a^{-1}b^{-3/2}.
\end{equation}
For cube-free $e=ab^2$, \cref{lem:rho} gives $\rho(p)\le d$ for $p\mid a$, $p\notin P_0$; $\rho(p^2)\le pd$ for $p\mid b$, $p\notin P_0$; and $\rho(p^{v_p(e)})\le p^2$ for $p\in P_0$. Hence $\rho(e)\le C_Pd^{\omega(e)}b$ with $C_P=\prod_{p\in P_0}p^2$, and
\begin{equation}\label{eq:w-rho}
|w(e)|\rho(e)\le C_P(2d)^{\omega(e)}a^{-1}b^{-1/2},\qquad\frac{|w(e)|\rho(e)}e\le C_P(2d)^{\omega(e)}a^{-2}b^{-5/2}.
\end{equation}

\emph{Three estimates.} Let $W\ge1$.
\begin{enumerate}[label=(\roman*)]
\item \emph{Tail.} For $e>W$ we have $1\le(e/W)^{1/2}$ and $e^{1/2}=a^{1/2}b$. By~\eqref{eq:w-rho},
\begin{equation}\label{eq:w-tail}
\sum_{e>W}\frac{|w(e)|\rho(e)}e\le W^{-1/2}C_P\sum_{a,b\ \mathrm{squarefree}}\frac{(2d)^{\omega(a)+\omega(b)}}{a^{3/2}b^{3/2}}=c_4W^{-1/2},
\end{equation}
with $c_4=C_P\prod_p\bigl(1+2dp^{-3/2}\bigr)^2$.
In particular $\sum_e|w(e)|\rho(e)/e<\infty$.
\item \emph{Short sum.} We have $(2d)^{\omega(n)}\le\tau(n)^{\log_2(2d)}\le C(f,\eps)n^\eps$ by \cref{lem:divisor}, where the constant depends on $d$ and $\eps$. By~\eqref{eq:w-rho}, using $\sum_{b\le B}b^{-1/2}\le2\sqrt B$,
\begin{equation}\label{eq:w-short}
\sum_{e\le W}|w(e)|\rho(e)\le C_PC(f,\eps)W^\eps\sum_{a\le W}\frac1a\sum_{b\le\sqrt W}b^{-1/2}\le2C_PC(f,\eps)W^{1/4+\eps}(1+\log W).
\end{equation}
\item \emph{Euler product.} The function $e\mapsto w(e)\rho(e)/e$ is multiplicative, and its series converges absolutely by~(i). By~\eqref{eq:euler} and the second expression for $\sigma_p$ in the proof of~(0),
\begin{equation}\label{eq:w-euler}
\sum_{e\ge1}\frac{w(e)\rho(e)}e=\prod_p\bigl(1+w(p)x_p+w(p^2)z_p\bigr)=\prod_p\frac{a_p+(b_p-a_p)x_p+(c_p-b_p)z_p}{a_p}=\frac1{A_h}\prod_p\sigma_p .
\end{equation}
\end{enumerate}

\emph{Truncation.} Let $y\ge1$ and $S=\sum_{k\le y}F(k)\hnat(f(k))$. For $1\le k\le y$ write $\hnat(f(k))=\sum_{e\mid f(k),\,e\le y}w(e)+T_k$ with $T_k=\sum_{e\mid f(k),\,e>y}w(e)$. By~\eqref{eq:w-bounds} and \cref{lem:divisor}, and since $2\le f(k)\le Y$ (as $f$ is increasing on $[1,\infty)$),
\[
|T_k|\le y^{-3/4}\sum_{e\mid f(k)}4^{\omega(e)}\le y^{-3/4}\tau(f(k))^3\le C(\eps)\,y^{-3/4}Y^\eps .
\]
By \cref{lem:shape}, $\sum_{k\le y}F(k)\le y\gamma_3F(y)\le\gamma_3\gamma_4Y^{3/2}$, so
\begin{equation}\label{eq:trunc}
\Bigl|\sum_{k\le y}F(k)T_k\Bigr|\le C(\eps)\gamma_3\gamma_4\,Y^{3/2+\eps}y^{-3/4}.
\end{equation}

\emph{Counting in residue classes.} For $e\in\N$ let $N(e)=\sum_{1\le k\le y,\ e\mid f(k)}F(k)$. This is a sum over the $\rho(e)$ residue classes $r\bmod e$ with $f(r)\equiv0\pmod e$. Intersect the at most $2d-1$ closed intervals of \cref{lem:shape} with $[1,y]$, and make them disjoint by removing from each its left endpoint whenever that point already belongs to the previous interval. This gives at most $2d-1$ disjoint intervals, closed or half-open, which cover $[1,y]$ and on each of which $F$ is monotone. Applying \cref{lem:monotone} on each of them, and summing,
\[
\Bigl|\sum_{\substack{1\le k\le y\\k\equiv r\,(e)}}F(k)-\frac1e\int_1^yF(t)\,dt\Bigr|\le3(2d-1)\sup_{[1,y]}F\le6d\gamma_3F(y).
\]
(Stating \cref{lem:monotone} for arbitrary intervals covers the closed first interval $[1,\beta]$ directly; no value of $F$ to the left of $1$ is needed.) Summing over the $\rho(e)$ classes,
\begin{equation}\label{eq:Ne}
N(e)=\frac{\rho(e)}e\cdot\frac23\bigl(Y^{3/2}-f(1)^{3/2}\bigr)+\theta_e\,6d\gamma_3\,\rho(e)F(y).
\end{equation}

\emph{Assembly.} By the truncation, $S=\sum_{e\le y}w(e)N(e)+\sum_{k\le y}F(k)T_k$. By~\eqref{eq:w-tail} and~\eqref{eq:w-euler} with $W=y$, $\sum_{e\le y}w(e)\rho(e)/e=A_h^{-1}\prod_p\sigma_p+\theta c_4y^{-1/2}$. Inserting~\eqref{eq:Ne} and using~\eqref{eq:w-short},~\eqref{eq:trunc} and $F(y)\le\gamma_4Y^{3/2}/y$,
\begin{align*}
S&=\frac23\bigl(Y^{3/2}-f(1)^{3/2}\bigr)\Bigl(\frac1{A_h}\prod_p\sigma_p+\theta c_4y^{-1/2}\Bigr)\\
&\qquad+\theta\,6d\gamma_3F(y)\sum_{e\le y}|w(e)|\rho(e)+\theta C(\eps)\gamma_3\gamma_4Y^{3/2+\eps}y^{-3/4}\\
&=\frac2{3A_h}\Bigl(\prod_p\sigma_p\Bigr)Y^{3/2}+O_{f,\eps}\Bigl(1+Y^{3/2}y^{-1/2}+Y^{3/2}y^{-3/4+\eps}(1+\log y)+Y^{3/2+\eps}y^{-3/4}\Bigr).
\end{align*}
The term $O_f(1)$ collects the contributions of $f(1)^{3/2}$. It is $\ll_fY^{3/2}y^{-1/2}$, because $Y^{3/2}y^{-1/2}\ge\gamma_1^{3/2}y^{(3d-1)/2}\ge\gamma_1^{3/2}$. For $\eps\le\frac18$ we have $y^{-3/4+\eps}(1+\log y)\ll_\eps y^{-1/2}$, and since $Y\ge2$ every error term is then $\ll_{f,\eps}Y^{3/2+\eps}y^{-1/2}$. For $\eps>\frac18$ the claim follows from the case $\eps=\frac18$. Multiplying by $IA_h$ gives
\[
\sum_{k\le y}F(k)h(f(k))=IA_h\,S=\frac23I\Bigl(\prod_p\sigma_p\Bigr)Y^{3/2}+O_{f,\eps}\bigl(Y^{3/2+\eps}y^{-1/2}\bigr)=C_fY^{3/2}+O_{f,\eps}\bigl(Y^{3/2+\eps}y^{-1/2}\bigr).
\]
This proves~(a). All implied constants depend only on $f$ and $\eps$, through $d$, $\gamma_1,\dots,\gamma_4$, $C_P$, $f(1)$ and the divisor-bound constants. \qed

\section{Proof of the main theorem}\label{sec:assembly}

Fix $f$ as in \cref{thm:main}. The finiteness and positivity of $C_f$ are \cref{prop:poly}(0). Let $x\ge2$ and $y\ge1$ be real with $Y=f(y)\le x$, and fix $\eps>0$.

\emph{Range of the moduli.} The hypothesis $f'(t)\ge1$ for all \emph{real} $t\ge1$ makes $f$ increasing on $[1,\infty)$. Hence every integer $1\le k\le y$ satisfies
\begin{equation}\label{eq:moduli-range}
2\le f(k)\le f(y)=Y\le x .
\end{equation}
So \cref{prop:single} applies to $q=f(k)$ with the same $x$, for every $k\le y$, with a constant independent of $k$. By~\eqref{eq:Vf}, $V_f(x,y)=\sum_{1\le k\le y}f'(k)V(x;f(k))$, and $f'(k)>0$.

\emph{Elementary inequalities.} Write $f(t)=\sum_{i\le d}\alpha_it^i$ and $\gamma_0=\sum_i|\alpha_i|$. For $y\ge1$, $Y\le\gamma_0y^d$, hence
\begin{equation}\label{eq:y-Y}
y^{-1/2}\le\gamma_0^{1/(2d)}\,Y^{-1/(2d)}.
\end{equation}
Since $Y\le x$, we have $Y^\eps\le x^\eps$ and $1+\log Y\le1+\log x\le C'(\eps)x^\eps$. With \cref{prop:poly}(b),(c) this gives
\begin{equation}\label{eq:sums}
\sum_{k\le y}\frac{f'(k)}{f(k)}\le c_8(f)C'(\eps)\,x^{\eps},\qquad\sum_{k\le y}f'(k)\le c_8(f)\,Y .
\end{equation}

\subsection{Part (ii)}

Apply \cref{prop:single}(i) with $\eps/2$ to each $q=f(k)$, multiply by $f'(k)$ and sum over $1\le k\le y$:
\[
\Bigl|V_f(x,y)-\sum_{k\le y}f'(k)\MT(x;f(k))\Bigr|\le c_7(\tfrac\eps2)\Bigl(x^{5/3+\eps/2}\sum_{k\le y}\frac{f'(k)}{f(k)}+\sqrt x\sum_{k\le y}f'(k)\Bigr).
\]
By~\eqref{eq:sums} with $\eps/2$, the right-hand side is at most
\[
c_7(\tfrac\eps2)\,c_8(f)\bigl(C'(\tfrac\eps2)\,x^{5/3+\eps}+x^{1/2}Y\bigr),
\]
which is part~(ii).

\subsection{Part (i)}

Apply \cref{prop:single}(ii) with $\eps/3$ to each $q=f(k)$ and use $f'(k)\sqrt{f(k)}=F(k)$:
\[
\Bigl|V_f(x,y)-\sqrt x\sum_{k\le y}F(k)h(f(k))\Bigr|\le c_7(\tfrac\eps3)\,(T_1+T_2+T_3),
\]
where the three sums are bounded as follows.
\begin{itemize}
\item $T_1=x^{1/3}\sum_{k\le y}f'(k)f(k)^{2/3+\eps/3}\le x^{1/3}Y^{2/3+\eps/3}\sum_{k\le y}f'(k)\le c_8(f)\,x^{1/3+\eps/3}Y^{5/3}$, by $f(k)\le Y\le x$ and~\eqref{eq:sums}.
\item $T_2=x^{5/3+\eps/3}\sum_{k\le y}f'(k)/f(k)\le c_8(f)C'(\tfrac\eps3)\,x^{5/3+2\eps/3}$, by~\eqref{eq:sums} with $\eps/3$.
\item $T_3=\sqrt x\sum_{k\le y}f'(k)\le c_8(f)\,x^{1/2}Y\le c_8(f)\,x^{1/2}Y^{3/2-1/(2d)}$, because $Y\ge2$ and $\frac12-\frac1{2d}\ge0$.
\end{itemize}
For the main term, \cref{prop:poly}(a) with $\eps/3$ and~\eqref{eq:y-Y} give
\begin{align*}
\Bigl|\sqrt x\sum_{k\le y}F(k)h(f(k))-C_fx^{1/2}Y^{3/2}\Bigr|&\le c_8(f,\tfrac\eps3)\gamma_0^{1/(2d)}\,x^{1/2}Y^{3/2+\eps/3-1/(2d)}\\
&\le c_8(f,\tfrac\eps3)\gamma_0^{1/(2d)}\,x^{1/2+\eps/3}Y^{3/2-1/(2d)}.
\end{align*}
By the triangle inequality,
\[
\bigl|V_f(x,y)-C_fx^{1/2}Y^{3/2}\bigr|\le c(f,\eps)\Bigl(x^{1/3+\eps}Y^{5/3}+x^{5/3+\eps}+x^{1/2+\eps}Y^{3/2-1/(2d)}\Bigr),
\]
with $c(f,\eps)=c_7(\tfrac\eps3)c_8(f)\bigl(2+C'(\tfrac\eps3)\bigr)+c_8(f,\tfrac\eps3)\gamma_0^{1/(2d)}$. This is part~(i).

\subsection{Part (iii)}

Fix $\delta\in(0,1/9)$. The interval $[x^{7/9+\delta},x^{1-\delta}]$ is nonempty for $x>1$ exactly because $7/9+\delta<1-\delta$. Assume $x^{7/9+\delta}\le Y\le x^{1-\delta}$, put $\eta=\min(\delta/12,1/(6d))$, and apply~(i) with $\eps=\min(\delta/12,1/(9d))$. We divide each error term by $x^{1/2}Y^{3/2}$:
\begin{enumerate}[label=\arabic*.]
\item $x^{-1/6+\eps}Y^{1/6}\le x^{-1/6+\eps+(1-\delta)/6}=x^{\eps-\delta/6}\le x^{-\delta/12}\le x^{-\eta}$, using $Y\le x^{1-\delta}$ and $\eps\le\delta/12$;
\item $x^{7/6+\eps}Y^{-3/2}\le x^{7/6+\eps-\frac32(7/9+\delta)}=x^{\eps-3\delta/2}\le x^{-\delta/12}\le x^{-\eta}$, using $Y\ge x^{7/9+\delta}$ and $\frac32\cdot\frac79=\frac76$;
\item $x^{\eps}Y^{-1/(2d)}\le x^{\eps-7/(18d)}\le x^{1/(9d)-7/(18d)}=x^{-5/(18d)}\le x^{-1/(6d)}\le x^{-\eta}$, using $Y\ge x^{7/9}$ and $\eps\le1/(9d)$.
\end{enumerate}
Hence $|V_f(x,y)-C_fx^{1/2}Y^{3/2}|\le3c(f,\eps)\,x^{-\eta}x^{1/2}Y^{3/2}$, where $\eps$ depends only on $\delta$ and $d$. As $0<C_f<\infty$, this is $V_f(x,y)=C_fx^{1/2}Y^{3/2}\bigl(1+O_{f,\delta}(x^{-\eta})\bigr)$. \qed

\begin{remark}
The exponent $\eta$ is conservative. The third ratio is in fact at most $x^{-5/(18d)}$, and the second at most $x^{-17\delta/12}$. We keep the exponent as it was stated in the archived proof.
\end{remark}

\section{Numerical evidence}\label{sec:numerics}

This section compares \cref{thm:main} with exact computations of $V_f(x,y)$. Everything here is evidence and plays no part in the proofs. All tables and plots are regenerated by \nolinkurl{manuscript/reproduce.py} (see \nolinkurl{manuscript/README.md}). That script re-runs the accepted computation from its source code and recompiles its C\texttt{++} part; it also computes the theoretical quantities. Numbers quoted in this section as bounds or ranges (``at most'', ``between'', ``lies in'') are rounded outwards, so that the printed statement holds for the computed values; numbers quoted as approximate values (``about'', ``approximately'', $\approx$) are rounded to the nearest digit.

\subsection{What was computed}\label{sec:num-method}

\emph{Grid.} The three polynomials are $f(k)=k^2+1$, $k^3+2$ and $2k+1$, with weights $f'(k)=2k$, $3k^2$ and $2$. All three satisfy~\eqref{eq:f-hyp}. The values of $x$ are $x=2.5\cdot10^5\cdot2^j$ for $0\le j\le10$, so the largest is $x=2.56\cdot10^8$. For each $x$ and each target ratio $r\in\{0.01,0.02,0.05,0.1,0.2,0.4,0.7,1\}$, $y$ is the largest integer with $f(y)\le rx$. At the two largest $x$ the ratio $r=0.005$ was added as a diagnostic. The quantity reported is the normalised variance
\[
N(x,r)=\frac{V_f(x,y)}{x^{1/2}f(y)^{3/2}},
\]
computed with the actual ratio $f(y)/x$. By \cref{thm:main}(iii), $N(x,r)\to C_f$ when $x\to\infty$ and $f(y)/x\to0$ suitably.

\emph{Method.} The computation (archived numerical snapshot, files \texttt{compute.py} and \texttt{counts.cpp}) is exact in all integer quantities.
\begin{enumerate}[label=(\arabic*)]
\item The squarefree indicator $\mu^2(n)$, $n\le x$, is sieved. All autocorrelations
\[
C(h)=\sum_{1\le n\le x-h}\mu^2(n)\mu^2(n+h),\qquad0\le h\le x,
\]
are obtained from one number-theoretic transform modulo the prime $p_0=3221225473$, with primitive root~$5$. Every $C(h)$ is at most $x<p_0$, so no information is lost.
\item As in the proof of \cref{prop:identity}, $\sum_lQ(x;q,l)^2=Q(x)+2\sum_{m\ge1}C(mq)$ for every modulus $q=f(k)$.
\item The densities are $g(q,l)=\frac6{\pi^2}\tilde g(q,l)$ with an explicit rational $\tilde g(q,l)$ equivalent to \cref{prop:densities}(i),(iii). The coefficients $\sum_lQ(x;q,l)\tilde g(q,l)$ and $\sum_l\tilde g(q,l)^2$ are computed as exact rationals, using the counts of squarefree multiples of $d$ as in the proof of \cref{prop:trivial}.
\item Finally $V(x;q)=\sum_lQ^2-2x\sum_lgQ+x^2\sum_lg^2$, as in~\eqref{eq:id1}, is accumulated with the weights $f'(k)$ in $70$-digit decimal arithmetic.
\end{enumerate}
The snapshot validates the implementation in several ways. It checks the transform against direct lag sums at every $x$. It compares $\tilde g$ with brute-force local densities for all $q\le50$ and with empirical frequencies at $x=10^6$. At $x=257$ it compares every generated modulus with direct residue counts. During the original run, the operator and an integrator also recomputed some rows independently by direct counting (Appendix~\ref{app:provenance}).

\emph{Re-run.} For this write-up the snapshot code was run again, unchanged, after its SHA-256 hashes had been checked against the acceptance record. The C\texttt{++} part was compiled locally with \texttt{clang++ -O3}. All $270$ rows ($90$ per polynomial, including the diagnostic rows) reproduce the archived integer counts exactly. The values of $V_f$ agree to all stored digits (relative differences at most $\RerunMaxRelV$), and all fitted quantities are identical. The full grid needs about $5$\,GB of memory; the running times of the re-run that produced the tables are recorded in \texttt{reproduction.json} (fields \texttt{snapshot\_rerun.seconds} and \texttt{snapshot\_rerun.total\_seconds\_of\_reproduce\_py}).

\subsection{The constants \texorpdfstring{$C_f$}{C\_f}}\label{sec:num-constants}

For the three polynomials the root counts are:
\begin{itemize}
\item $k^2+1$: $\rho(2)=1$, $\rho(4)=0$, and $\rho(p)=\rho(p^2)=1+\bigl(\frac{-1}{p}\bigr)$ for odd $p$;
\item $k^3+2$: $\rho(2)=\rho(3)=1$, $\rho(4)=\rho(9)=0$; for $p\ge5$, $\rho(p^2)=\rho(p)$, with $\rho(p)=1$ if $p\equiv2\pmod 3$, and, if $p\equiv1\pmod3$, $\rho(p)=3$ or $0$ according as $(-2)^{(p-1)/3}\equiv1\pmod p$ or not;
\item $2k+1$: $\rho(2)=\rho(4)=0$, and $\rho(p)=\rho(p^2)=1$ for odd $p$.
\end{itemize}
For primes not dividing the discriminant all roots are simple, and each root modulo $p$ lifts to exactly one root modulo $p^2$, which gives $\rho(p^2)=\rho(p)$. The formulas were checked by brute force modulo $p$ and $p^2$ for all $p<\RhoBound$. None of the three polynomials has a content prime, so $P_0=\emptyset$, and~\eqref{eq:sigma-close} holds for every $p$ with $K=3+2d$. Hence, for $N\ge30$,
\[
\Bigl|\log\prod_{p>N}\sigma_p\Bigr|\le\frac{K}{1-K/N^2}\sum_{p>N}\frac1{p^2}\le\frac{8K}{1-K/N^2}\Bigl(\frac1{N^2}+\frac1{30N}\Bigr).
\]
Here we used that every prime $p>N\ge30$ lies in one of the $8$ reduced residue classes modulo $30$, and that each class contributes at most $N^{-2}+(30N)^{-1}$. \Cref{tab:constants} lists the partial values $\frac23I\prod_{p\le N}\sigma_p$ and the resulting interval for $C_f$ with $N=10^7$. The partial values are rounded to nine decimals; the interval endpoints are rounded outwards, so that each printed interval contains the computed one. The intervals account rigorously for the omitted primes. They do not account for floating-point rounding in the partial products, which are computed in double precision. As a check of that rounding (not a certified bound), the partial products for $N=10^5$ were recomputed in $40$-digit arithmetic; they agree with the double-precision values to a relative difference of at most $\FloatCheck$.

\begin{table}[htbp]
\centering
\footnotesize\setlength{\tabcolsep}{3.5pt}
\begin{tabular}{lrrrrrl}
\toprule
$f$ & $N=10^{3}$ & $N=10^{4}$ & $N=10^{5}$ & $N=10^{6}$ & $N=10^{7}$ & interval for $C_f$ \\
\midrule
$k^2+1$ & 0.138152461 & 0.138088027 & 0.138083068 & 0.138082663 & 0.138082629 & $[0.138082603,\,0.138082655]$ \\
$k^3+2$ & 0.135675596 & 0.135612820 & 0.135607954 & 0.135607555 & 0.135607522 & $[0.135607489,\,0.135607555]$ \\
$2k+1$ & 0.142685343 & 0.142618897 & 0.142613765 & 0.142613346 & 0.142613311 & $[0.142613291,\,0.142613330]$ \\
$t$ (formal) & 0.128543562 & 0.128483702 & 0.128479079 & 0.128478702 & 0.128478670 & $[0.128478652,\,0.128478687]$ \\
Vaughan $c_2$ & 0.128510919 & 0.128481180 & 0.128478873 & 0.128478684 & 0.128478668 & \\
\bottomrule
\end{tabular}

\caption{Partial Euler products $\frac23I\prod_{p\le N}\sigma_p$ and intervals for $C_f$. The row ``$t$ (formal)'' uses $\rho\equiv1$. The last row is Vaughan's constant~\eqref{eq:vaughan-c2} with $C_1$ truncated at $p\le N$; its limit equals that of the row above (\cref{rem:vaughan}).}
\label{tab:constants}
\end{table}

Apart from floating-point rounding, the intervals therefore fix the digits
\[
C_{k^2+1}=\CfQuadCert\ldots,\qquad C_{k^3+2}=\CfCubicCert\ldots,\qquad C_{2k+1}=\CfLinCert\ldots .
\]
The partial products at $N=10^7$ give the approximate values $C_{k^2+1}\approx\CfQuad$, $C_{k^3+2}\approx\CfCubic$ and $C_{2k+1}\approx\CfLin$; the digits beyond those fixed above are only indicative. For $f(t)=t$ the formal value $\frac23I\prod_{p\le 10^7}\sigma_p\approx\CfT$ agrees with Vaughan's $c_2\approx\VaughanCtwo$ (with $C_1$ truncated at the same point), as \cref{rem:vaughan} predicts. The local identity behind that remark, the identities of \cref{lem:local}, and $\zeta(-\frac12)=-\zeta(\frac32)/(4\pi)$ were also verified with computer algebra and $80$-digit arithmetic (\texttt{reproduce.py}, function \texttt{symbolic\_local\_identity}).

\subsection{The data and the predicted constant}\label{sec:num-data}

\Cref{tab:largest} shows $N(x,r)$ at the largest $x$, together with the relative change from $x=1.28\cdot10^8$. The full tables are in Appendix~\ref{app:tables}. Two features are visible.

First, at fixed $r$ the normalised values are stable in $x$. Between the two largest sizes they change by at most $\StabCubic\%$ for $k^3+2$, and by at most $\StabOther\%$ for $k^2+1$ and $2k+1$ (see also \cref{fig:x}). Fits of $\log V_f$ against $\log x$ at fixed target ratio give $\alpha+\beta\approx2$ (\cref{tab:fits}), consistent with the homogeneous main term $x^{1/2}Y^{3/2}$.

Second, $N(x,r)$ depends clearly on $r$. It rises from values between $\NDiagMin$ and $\NDiagMax$ at $r=0.005$ to values between $\NOneMin$ and $\NOneMax$ at $r=1$. A fit of $\log V_f$ against $\log f(y)$ at fixed $x$ therefore gives exponents $\beta$ between $\BetaMin$ and $\BetaMax$ instead of $\frac32$ (\cref{tab:fits}). This does not contradict \cref{thm:main}. Part (iii) concerns $f(y)/x\to0$, and at a fixed ratio $r$ the error term $x^{1/3+\eps}Y^{5/3}$ of part (i) is as large as the main term. The ratio dependence is captured by the exact-shape statement~(ii); see \cref{sec:num-shape}.

\begin{table}[htbp]
\centering
\footnotesize\setlength{\tabcolsep}{3.5pt}
\begin{tabular}{lrrrrrrrrr}
\toprule
$r$ & $0.005$ & $0.01$ & $0.02$ & $0.05$ & $0.1$ & $0.2$ & $0.4$ & $0.7$ & $1$ \\
\midrule
$k^2+1$ & 0.14039 & 0.14086 & 0.14292 & 0.14446 & 0.15098 & 0.15655 & 0.16278 & 0.16993 & 0.17110 \\
\quad change (\%) & $-0.0233$ & $-0.0363$ & $+0.0051$ & $-0.0102$ & $-0.0118$ & $-0.0043$ & $-0.0083$ & $-0.0003$ & $-0.0030$ \\
$k^3+2$ & 0.14076 & 0.13999 & 0.14120 & 0.14391 & 0.14848 & 0.15234 & 0.15550 & 0.16399 & 0.16671 \\
\quad change (\%) & $+0.4079$ & $-0.9050$ & $-0.5298$ & $-0.3279$ & $-0.3250$ & $-0.0824$ & $-0.2807$ & $-0.1486$ & $-0.1688$ \\
$2k+1$ & 0.14484 & 0.14460 & 0.14837 & 0.14934 & 0.15670 & 0.16171 & 0.16245 & 0.18239 & 0.18369 \\
\quad change (\%) & $+0.0003$ & $-0.0000$ & $+0.0002$ & $+0.0005$ & $+0.0002$ & $+0.0001$ & $+0.0001$ & $-0.0000$ & $-0.0000$ \\
\bottomrule
\end{tabular}

\caption{$N(x,r)=V_f(x,y)/(x^{1/2}f(y)^{3/2})$ at $x=2.56\cdot10^8$ for the target ratios $r$ (the column $0.005$ is the diagnostic ratio). The ``change'' rows give the relative change, in percent, from $x=1.28\cdot10^8$ to $x=2.56\cdot10^8$.}
\label{tab:largest}
\end{table}

\emph{Extrapolation to $r\to0$.} The snapshot estimates $\lim_{r\to0}N$ empirically. At the largest $x$ it fits the model $N(r)=C+a\sqrt r+br$ by unweighted least squares to the rows with $r$ at or above a ``stable'' cutoff. The cutoff is the smallest requested ratio such that every row at or above it changes by less than $1\%$ between the two largest $x$; here it is $\StableCutoff$ for all three polynomials. The reported uncertainty is a \emph{heuristic sensitivity envelope}: the largest of twice the formal standard error of the intercept, the spread over neighbouring cutoffs, the shift to the previous $x$, the shifts under the alternative models $C+a\sqrt r$ and $C+a\sqrt r+br+cr^{3/2}$, and the shifts when the fitted window is cut at $r\le0.1$, $0.2$ or $0.4$. It is not a confidence interval and not an error bound, and the model $C+a\sqrt r+br$ itself is not derived from the theorem. \Cref{tab:fits} compares the extrapolated constants with $C_f$. The relative differences are approximately $\RelQuad\%$, $\RelCubic\%$ and $\RelLin\%$ for $k^2+1$, $k^3+2$ and $2k+1$, all within the envelopes. \Cref{tab:sensitivity} shows how strongly the extrapolated value depends on the model. For $k^2+1$ the alternatives range from about $\SensQuadMin$ to $\SensQuadMax$ and bracket $C_f$, so the difference of $\RelQuad\%$ reflects the choice of model rather than the constant. \Cref{fig:ratio} shows the data, the fitted curves and $C_f$.

\begin{table}[htbp]
\centering
\small
\begin{tabular}{lrrrrrr}
\toprule
$f$ & $\alpha+\beta$ (fixed $r$) & $\beta$ (fixed $x$) & $C_{\mathrm{extr}}\pm$ envelope & $C_f$ & $\frac{C_{\mathrm{extr}}-C_f}{C_f}$ & $\frac{C_{\mathrm{extr}}-C_f}{\mathrm{envelope}}$ \\
\midrule
$k^2+1$ & $2.0008\pm0.0006$ & $1.545\pm0.004$ & $0.1345\pm0.0073$ & $0.13808$ & $-2.61\%$ & $-0.50$ \\
$k^3+2$ & $1.9978\pm0.0021$ & $1.538\pm0.004$ & $0.1363\pm0.0024$ & $0.13561$ & $+0.53\%$ & $+0.30$ \\
$2k+1$ & $2.0001\pm0.0001$ & $1.551\pm0.007$ & $0.1417\pm0.0082$ & $0.14261$ & $-0.65\%$ & $-0.11$ \\
\bottomrule
\end{tabular}

\caption{Empirical exponents and extrapolated constants from the snapshot fits, compared with $C_f$ from \cref{tab:constants}. ``$\alpha+\beta$'' is the mean of the eight fixed-ratio slopes of $\log V_f$ against $\log x$ (with the across-ratio standard error); ``$\beta$'' is the fixed-$x$ slope against $\log f(y)$ at the largest $x$ (with its regression standard error). The last column is $(C_{\mathrm{extr}}-C_f)$ divided by the heuristic envelope; a value of absolute size at most $1$ means that $C_f$ lies within the envelope.}
\label{tab:fits}
\end{table}

\begin{table}[htbp]
\centering
\footnotesize
\begin{tabular}{lrrrrrr}
\toprule
$f$ & primary & $C+a\sqrt r$ & $C+a\sqrt r+br+cr^{3/2}$ & $r\le 0.1$ & cutoff $0.005$ & previous $x$ \\
\midrule
$k^2+1$ & 0.13447 & 0.13837 & 0.13714 & 0.14176 & 0.13547 & 0.13451 \\
$k^3+2$ & 0.13632 & 0.13759 & 0.13569 & 0.13856 & 0.13728 & 0.13763 \\
$2k+1$ & 0.14169 & 0.14086 & 0.14145 & 0.14416 & 0.14188 & 0.14169 \\
\bottomrule
\end{tabular}

\caption{Sensitivity of the extrapolated constant: the primary fit, the two alternative models, the window $r\le0.1$, the lower cutoff $0.005$, and the primary fit at the previous size $x=1.28\cdot10^8$.}
\label{tab:sensitivity}
\end{table}

\begin{figure}[htbp]
\centering
\begin{tikzpicture}
\begin{axis}[width=0.82\textwidth,height=0.5\textwidth,xlabel={$\sqrt{r}$, with $r=f(y)/x$},ylabel={$N(x,r)$},xmin=-0.02,xmax=1.02,ymin=0.13,ymax=0.19,
  legend pos=north west,legend cell align=left,legend style={font=\footnotesize},tick label style={font=\footnotesize},
  yticklabel style={/pgf/number format/fixed,/pgf/number format/precision=3}]
\addplot[only marks,mark=*,mark size=1.6pt,blue] table[x=sqrtr,y=N]{generated/fig_ratio_0.dat};
\addlegendentry{$k^2+1$}
\addplot[only marks,mark=square*,mark size=1.6pt,red] table[x=sqrtr,y=N]{generated/fig_ratio_1.dat};
\addlegendentry{$k^3+2$}
\addplot[only marks,mark=triangle*,mark size=2pt,black!60!green] table[x=sqrtr,y=N]{generated/fig_ratio_2.dat};
\addlegendentry{$2k+1$}
\addplot[blue,dashed,forget plot] table[x=sqrtr,y=N]{generated/fig_fit_0.dat};
\addplot[red,dashed,forget plot] table[x=sqrtr,y=N]{generated/fig_fit_1.dat};
\addplot[black!60!green,dashed,forget plot] table[x=sqrtr,y=N]{generated/fig_fit_2.dat};
\addplot[only marks,mark=o,mark size=3pt,blue,thick,forget plot] table[x=sqrtr,y=N]{generated/fig_cf_0.dat};
\addplot[only marks,mark=o,mark size=3pt,red,thick,forget plot] table[x=sqrtr,y=N]{generated/fig_cf_1.dat};
\addplot[only marks,mark=o,mark size=3pt,black!60!green,thick,forget plot] table[x=sqrtr,y=N]{generated/fig_cf_2.dat};
\end{axis}
\end{tikzpicture}
\caption{Normalised variance at $x=2.56\cdot10^8$ against $\sqrt r$ (filled marks), the snapshot's empirical fits $C+a\sqrt r+br$ (dashed), and the proved constants $C_f$ (open circles at $r=0$).}
\label{fig:ratio}
\end{figure}

\begin{figure}[htbp]
\centering
\begin{tikzpicture}
\begin{axis}[width=0.6\textwidth,height=0.42\textwidth,xmode=log,xlabel={$x$},ylabel={$N(x,r)$},ymin=0.135,ymax=0.19,
  legend pos=outer north east,legend cell align=left,legend style={font=\footnotesize},tick label style={font=\footnotesize},
  yticklabel style={/pgf/number format/fixed,/pgf/number format/precision=3}]
\addplot[blue,mark=*,mark size=1.3pt] table[x=x,y=N]{generated/fig_x_0a.dat};\addlegendentry{$k^2+1$, $r=0.01$}
\addplot[blue,mark=o,mark size=1.3pt,dashed] table[x=x,y=N]{generated/fig_x_0b.dat};\addlegendentry{$k^2+1$, $r=1$}
\addplot[red,mark=square*,mark size=1.3pt] table[x=x,y=N]{generated/fig_x_1a.dat};\addlegendentry{$k^3+2$, $r=0.01$}
\addplot[red,mark=square,mark size=1.3pt,dashed] table[x=x,y=N]{generated/fig_x_1b.dat};\addlegendentry{$k^3+2$, $r=1$}
\addplot[black!60!green,mark=triangle*,mark size=1.6pt] table[x=x,y=N]{generated/fig_x_2a.dat};\addlegendentry{$2k+1$, $r=0.01$}
\addplot[black!60!green,mark=triangle,mark size=1.6pt,dashed] table[x=x,y=N]{generated/fig_x_2b.dat};\addlegendentry{$2k+1$, $r=1$}
\end{axis}
\end{tikzpicture}
\caption{Stability in $x$: $N(x,r)$ against $x$ for $r=0.01$ and $r=1$.}
\label{fig:x}
\end{figure}

\subsection{The exact-shape statement}\label{sec:num-shape}

\Cref{thm:main}(ii) replaces $C_fx^{1/2}Y^{3/2}$ by $\Sigma_{\MT}=\sum_{k\le y}f'(k)\MT(x;f(k))$. We evaluated $\Sigma_{\MT}$ for the rows listed in \cref{tab:shape}, together with $\sum_{k\le y}F(k)h(f(k))$, the main term of \cref{prop:single}(ii) summed over the moduli. We report
\[
N_{\MT}=\frac{\Sigma_{\MT}}{x^{1/2}Y^{3/2}},\qquad N_h=\frac{1}{Y^{3/2}}\sum_{k\le y}F(k)h(f(k)).
\]
$\MT(x;q)$ is computed from the factorisation $t=t_0t_1t_2$ of the proof of \cref{prop:MT}. For $t_0\le\sqrt{Z_1}$ the terms are summed directly. For $t_0>\sqrt{Z_1}$ we use $B(u)=u-u^2$ ($u<1$) and the closed forms
\[
A_0\sum_{(t,q)=1}\eta(t)=\frac6{\pi^2}\prod_{p\mid q}\frac{p^2-2}{p^2-1},\qquad A_0\sum_{(t,q)=1}\frac{\eta(t)}{t^2}=\frac{36}{\pi^4}\prod_{p\mid q}\frac{p^2(p^2-2)}{(p^2-1)^2},
\]
which follow from~\eqref{eq:euler} as in \cref{lem:eta}. The first $64$ terms of the subtracted partial sums are evaluated in $80$-digit arithmetic; $A_0$ and $\prod_pa_p$ are evaluated to high precision through the prime zeta function. The infinite series~\eqref{eq:MT} is thus evaluated without truncation, up to floating-point rounding. As an independent check, \cref{prop:identity} was verified with this implementation of $\MT$ against brute-force values of every other term, for $x=3000$ and all $q\le40$ (relative discrepancy at most $\SelfTest$). For $2k+1$ the number of moduli is about $x/2$, so only $x\le10^6$ was treated.

\begin{table}[htbp]
\centering
\small
\begin{tabular}{llrrrrrr}
\toprule
$f$ & $x$ & $r$ & $y$ & $N_{\mathrm{data}}$ & $N_{\mathrm{MT}}$ & $V_f/\Sigma_{\mathrm{MT}}$ & $N_h$ \\
\midrule
$k^2+1$ & $10^{6}$ & $0.01$ & 99 & 0.14185 & 0.14171 & 1.00100 & 0.13966 \\
 &  & $0.1$ & 316 & 0.15191 & 0.15191 & 0.99999 & 0.13898 \\
 &  & $1$ & 999 & 0.17136 & 0.17136 & 1.00003 & 0.13825 \\
 & $1.6\cdot 10^{7}$ & $0.01$ & 399 & 0.14117 & 0.14118 & 0.99996 & 0.13853 \\
 &  & $0.1$ & 1264 & 0.15115 & 0.15115 & 1.00003 & 0.13825 \\
 &  & $1$ & 3999 & 0.17115 & 0.17115 & 1.00000 & 0.13812 \\
 & $2.56\cdot 10^{8}$ & $0.005$ & 1131 & 0.14039 & 0.14039 & 0.99998 & 0.13829 \\
 &  & $0.01$ & 1599 & 0.14086 & 0.14086 & 0.99999 & 0.13820 \\
 &  & $0.1$ & 5059 & 0.15098 & 0.15097 & 1.00001 & 0.13812 \\
 &  & $1$ & 15999 & 0.17110 & 0.17110 & 1.00000 & 0.13809 \\
\midrule
$k^3+2$ & $10^{6}$ & $0.01$ & 21 & 0.15356 & 0.15361 & 0.99968 & 0.15544 \\
 &  & $0.1$ & 46 & 0.15503 & 0.15499 & 1.00023 & 0.14231 \\
 &  & $1$ & 99 & 0.17071 & 0.17071 & 1.00003 & 0.13890 \\
 & $1.6\cdot 10^{7}$ & $0.01$ & 54 & 0.14240 & 0.14236 & 1.00024 & 0.14132 \\
 &  & $0.1$ & 116 & 0.14975 & 0.14974 & 1.00009 & 0.13844 \\
 &  & $1$ & 251 & 0.16787 & 0.16787 & 1.00000 & 0.13708 \\
 & $2.56\cdot 10^{8}$ & $0.005$ & 108 & 0.14076 & 0.14079 & 0.99982 & 0.13851 \\
 &  & $0.01$ & 136 & 0.13999 & 0.14000 & 0.99995 & 0.13776 \\
 &  & $0.1$ & 294 & 0.14848 & 0.14848 & 1.00000 & 0.13673 \\
 &  & $1$ & 634 & 0.16671 & 0.16671 & 1.00000 & 0.13602 \\
\midrule
$2k+1$ & $2.5\cdot 10^{5}$ & $0.01$ & 1249 & 0.14376 & 0.14466 & 0.99381 & 0.14268 \\
 &  & $0.1$ & 12499 & 0.15670 & 0.15671 & 0.99992 & 0.14262 \\
 &  & $1$ & 124999 & 0.18370 & 0.18369 & 1.00002 & 0.14261 \\
 & $10^{6}$ & $0.01$ & 4999 & 0.14439 & 0.14461 & 0.99848 & 0.14263 \\
 &  & $0.1$ & 49999 & 0.15670 & 0.15670 & 0.99997 & 0.14262 \\
 &  & $1$ & 499999 & 0.18369 & 0.18369 & 0.99999 & 0.14261 \\
\bottomrule
\end{tabular}

\caption{Exact-shape comparison: the data $N_{\mathrm{data}}=N(x,r)$, the normalised sum $N_{\MT}$ of the exact main terms, the ratio $V_f/\Sigma_{\MT}$, and the normalised $h$-average $N_h$. Rows with $r\in\{0.005,0.01,0.1,1\}$ are shown; \texttt{reproduction.json} contains all $\ShapeRows$ computed rows.}
\label{tab:shape}
\end{table}

Over all $\ShapeRows$ computed rows, $V_f/\Sigma_{\MT}$ lies in $[\ShapeMin,\ShapeMax]$. At $x=2.56\cdot10^8$ it differs from $1$ by at most $\ShapeDevLargest$. So the exact main term of \cref{thm:main}(ii) accounts for the whole ratio dependence seen in \cref{tab:largest}. This agreement is much better than the proved error term $x^{5/3+\eps}+x^{1/2}Y$. Relative to $x^{1/2}Y^{3/2}$, and ignoring the unknown constant and the factor $x^\eps$, that error term is of size about $x^{-1/3}r^{-3/2}$. At $x=2.56\cdot10^8$ this is about $\XLargeInvCubeSmall$ for $r=0.01$ and $\XLargeInvCubeOne$ for $r=1$. The pair-count errors $D_x(hq)$ evidently cancel substantially when summed over $h$, which the proof does not exploit. The quantity $N_h$ approaches $C_f$ as $y$ grows, as \cref{prop:poly}(a) predicts. For $k^3+2$ the convergence is slower and less regular, because $y$ is only a few hundred and the error is of order $y^{-1/2}$. The difference $N_{\MT}-N_h$ is the secondary term of the sawtooth sums. It is small for small $r$ and grows with $r$. This is consistent with the relative bound $(Y/x)^{1/6}$ (up to $x^{\eps}$) that \cref{prop:MT} implies, but in the data the difference behaves roughly like a multiple of $\sqrt r$ (\cref{fig:shape}).

\begin{figure}[htbp]
\centering
\begin{tikzpicture}
\begin{groupplot}[group style={group size=2 by 1,horizontal sep=1.6cm},width=0.47\textwidth,height=0.4\textwidth,
  xlabel={$\sqrt r$},xmin=-0.02,xmax=1.02,ymin=0.13,ymax=0.175,tick label style={font=\footnotesize},
  yticklabel style={/pgf/number format/fixed,/pgf/number format/precision=3},legend style={font=\footnotesize},legend cell align=left]
\nextgroupplot[title={$k^2+1$, $x=2.56\cdot10^8$},ylabel={normalised value},legend pos=north west]
\addplot[only marks,mark=*,mark size=1.6pt,blue] table[x=sqrtr,y=Ndata]{generated/fig_shape_0.dat};\addlegendentry{$N_{\mathrm{data}}$}
\addplot[blue,thin] table[x=sqrtr,y=NMT]{generated/fig_shape_0.dat};\addlegendentry{$N_{\MT}$}
\addplot[black,dashed,mark=x] table[x=sqrtr,y=Nh]{generated/fig_shape_0.dat};\addlegendentry{$N_h$}
\addplot[only marks,mark=o,mark size=3pt,thick,black] table[x=sqrtr,y=N]{generated/fig_cf_0.dat};\addlegendentry{$C_f$}
\nextgroupplot[title={$k^3+2$, $x=2.56\cdot10^8$}]
\addplot[only marks,mark=square*,mark size=1.6pt,red] table[x=sqrtr,y=Ndata]{generated/fig_shape_1.dat};
\addplot[red,thin] table[x=sqrtr,y=NMT]{generated/fig_shape_1.dat};
\addplot[black,dashed,mark=x] table[x=sqrtr,y=Nh]{generated/fig_shape_1.dat};
\addplot[only marks,mark=o,mark size=3pt,thick,black] table[x=sqrtr,y=N]{generated/fig_cf_1.dat};
\end{groupplot}
\end{tikzpicture}
\caption{Data, exact main term $N_{\MT}$, averaged asymptotic main term $N_h$ and $C_f$ at the largest $x$. Both panels use the legend of the left panel: filled marks are the data, the solid line is $N_{\MT}$, the dashed line with crosses is $N_h$, and the open circle at $r=0$ is $C_f$.}
\label{fig:shape}
\end{figure}

\FloatBarrier
\subsection{Limitations}\label{sec:num-limits}

\begin{itemize}
\item The computations stop at $x=2.56\cdot10^8$. With the actual ratio $\varrho=f(y)/x$, a row lies in the range $x^{7/9+\delta}\le f(y)\le x^{1-\delta}$ of \cref{thm:main}(iii) exactly when
\[
\delta\le\delta_{\mathrm{row}}:=\min\Bigl(\frac{\log\varrho}{\log x}+\frac29,\ -\frac{\log\varrho}{\log x}\Bigr).
\]
At $x=2.56\cdot10^8$ we have $x^{-2/9}\approx\XTwoNinths$. The rows with $r\le0.01$ have $\varrho\le r<x^{-2/9}$, so $\delta_{\mathrm{row}}<0$ and they lie outside the range for every $\delta>0$. The rows with $r=1$ have $\varrho$ close to $1$ and $\delta_{\mathrm{row}}\le\DeltaRowOneMax$, so they lie in the range only for very small $\delta$; the implied constant in~(iii) depends on $\delta$ and is not explicit. Over the whole grid at this $x$, $\delta_{\mathrm{row}}$ is largest, about $\DeltaRowMax$, at $r=\DeltaRowMaxR$. For each fixed $\delta>0$, the rows with a fixed target ratio $r$ leave the range as $x$ grows, since they need $x\le\varrho^{-1/\delta}$, where $\varrho\approx r$. The values of $\delta_{\mathrm{row}}$ for all rows at the largest $x$ are recorded in \texttt{reproduction.json}.
\item The implied constants in \cref{thm:main} are not explicit, so no finite computation can confirm or refute it. The data can only be checked for consistency.
\item The extrapolated constants depend on an empirical model, and their uncertainties are heuristic envelopes, not confidence intervals. The agreement in \cref{tab:fits} is therefore qualitative evidence for the value of $C_f$, and the theorem does not rest on it.
\item The values of $C_f$ in \cref{tab:constants} come with a rigorous bound for the omitted primes, but not with a certified bound for floating-point rounding; the $40$-digit recomputation for $N=10^5$ is only a check.
\item The exact-shape comparison uses double-precision arithmetic, with $80$-digit evaluation only where cancellation occurs. Its agreement with the identity self-test is the main check of the implementation. For $2k+1$ it covers only $x\le10^6$.
\end{itemize}


\clearpage
\bibliographystyle{plainurl}
\bibliography{references}

\appendix
\section{Full tables of normalised values}\label{app:tables}

The tables give $N(x,r)=V_f(x,y)/(x^{1/2}f(y)^{3/2})$ for all sizes and target ratios of the snapshot grid, regenerated by \texttt{reproduce.py}. The integer $y$ is the largest with $f(y)\le rx$, and the normalisation uses the actual value $f(y)$.

\begin{table}[ht]
\centering
\small
\begin{tabular}{lrrrrrrrr}
\toprule
$x$ & $0.01$ & $0.02$ & $0.05$ & $0.1$ & $0.2$ & $0.4$ & $0.7$ & $1$ \\
\midrule
$2.5\cdot 10^{5}$ & 0.142649 & 0.146052 & 0.146549 & 0.152224 & 0.157416 & 0.163777 & 0.170589 & 0.171642 \\
$5\cdot 10^{5}$ & 0.142565 & 0.143625 & 0.145049 & 0.151489 & 0.157419 & 0.163281 & 0.170311 & 0.171372 \\
$10^{6}$ & 0.141855 & 0.143162 & 0.144886 & 0.151912 & 0.156975 & 0.163179 & 0.170250 & 0.171362 \\
$2\cdot 10^{6}$ & 0.141895 & 0.143860 & 0.145330 & 0.151304 & 0.156884 & 0.163051 & 0.170098 & 0.171269 \\
$4\cdot 10^{6}$ & 0.141686 & 0.143302 & 0.144687 & 0.151276 & 0.156752 & 0.162938 & 0.170032 & 0.171220 \\
$8\cdot 10^{6}$ & 0.141162 & 0.143155 & 0.144704 & 0.151199 & 0.156712 & 0.162890 & 0.170034 & 0.171183 \\
$1.6\cdot 10^{7}$ & 0.141175 & 0.143090 & 0.144622 & 0.151153 & 0.156632 & 0.162835 & 0.169986 & 0.171153 \\
$3.2\cdot 10^{7}$ & 0.141046 & 0.143021 & 0.144584 & 0.151056 & 0.156601 & 0.162812 & 0.169957 & 0.171141 \\
$6.4\cdot 10^{7}$ & 0.140949 & 0.142966 & 0.144505 & 0.151012 & 0.156578 & 0.162792 & 0.169935 & 0.171120 \\
$1.28\cdot 10^{8}$ & 0.140909 & 0.142912 & 0.144472 & 0.150993 & 0.156556 & 0.162789 & 0.169926 & 0.171106 \\
$2.56\cdot 10^{8}$ & 0.140858 & 0.142919 & 0.144457 & 0.150976 & 0.156550 & 0.162775 & 0.169925 & 0.171101 \\
\bottomrule
\end{tabular}

\caption{$f(k)=k^2+1$.}
\end{table}

\begin{table}[ht]
\centering
\small
\begin{tabular}{lrrrrrrrr}
\toprule
$x$ & $0.01$ & $0.02$ & $0.05$ & $0.1$ & $0.2$ & $0.4$ & $0.7$ & $1$ \\
\midrule
$2.5\cdot 10^{5}$ & 0.159913 & 0.160232 & 0.160971 & 0.159910 & 0.163172 & 0.161392 & 0.169394 & 0.172387 \\
$5\cdot 10^{5}$ & 0.155112 & 0.162671 & 0.156355 & 0.161317 & 0.159404 & 0.159028 & 0.167785 & 0.170862 \\
$10^{6}$ & 0.153562 & 0.149407 & 0.154182 & 0.155026 & 0.155087 & 0.158578 & 0.166278 & 0.170713 \\
$2\cdot 10^{6}$ & 0.148041 & 0.147909 & 0.151134 & 0.150024 & 0.154648 & 0.157691 & 0.166923 & 0.169599 \\
$4\cdot 10^{6}$ & 0.144716 & 0.147763 & 0.146250 & 0.151598 & 0.154780 & 0.157141 & 0.165214 & 0.168422 \\
$8\cdot 10^{6}$ & 0.141958 & 0.144833 & 0.145867 & 0.149564 & 0.154211 & 0.157510 & 0.165585 & 0.168245 \\
$1.6\cdot 10^{7}$ & 0.142399 & 0.144360 & 0.144416 & 0.149753 & 0.153584 & 0.156699 & 0.164844 & 0.167875 \\
$3.2\cdot 10^{7}$ & 0.142149 & 0.142778 & 0.144311 & 0.149048 & 0.153051 & 0.156096 & 0.164601 & 0.167379 \\
$6.4\cdot 10^{7}$ & 0.140427 & 0.142955 & 0.143925 & 0.149586 & 0.152600 & 0.155833 & 0.164525 & 0.167244 \\
$1.28\cdot 10^{8}$ & 0.141268 & 0.141954 & 0.144385 & 0.148961 & 0.152464 & 0.155936 & 0.164234 & 0.166987 \\
$2.56\cdot 10^{8}$ & 0.139990 & 0.141202 & 0.143912 & 0.148477 & 0.152339 & 0.155499 & 0.163990 & 0.166706 \\
\bottomrule
\end{tabular}

\caption{$f(k)=k^3+2$.}
\end{table}

\begin{table}[ht]
\centering
\small
\begin{tabular}{lrrrrrrrr}
\toprule
$x$ & $0.01$ & $0.02$ & $0.05$ & $0.1$ & $0.2$ & $0.4$ & $0.7$ & $1$ \\
\midrule
$2.5\cdot 10^{5}$ & 0.143761 & 0.148093 & 0.149307 & 0.156698 & 0.161734 & 0.162453 & 0.182396 & 0.183696 \\
$5\cdot 10^{5}$ & 0.144153 & 0.148232 & 0.149327 & 0.156723 & 0.161716 & 0.162465 & 0.182393 & 0.183694 \\
$10^{6}$ & 0.144392 & 0.148302 & 0.149320 & 0.156700 & 0.161696 & 0.162452 & 0.182388 & 0.183690 \\
$2\cdot 10^{6}$ & 0.144437 & 0.148283 & 0.149305 & 0.156686 & 0.161701 & 0.162446 & 0.182389 & 0.183691 \\
$4\cdot 10^{6}$ & 0.144507 & 0.148327 & 0.149321 & 0.156694 & 0.161706 & 0.162449 & 0.182390 & 0.183693 \\
$8\cdot 10^{6}$ & 0.144593 & 0.148372 & 0.149349 & 0.156706 & 0.161713 & 0.162452 & 0.182390 & 0.183692 \\
$1.6\cdot 10^{7}$ & 0.144609 & 0.148375 & 0.149350 & 0.156706 & 0.161712 & 0.162452 & 0.182390 & 0.183691 \\
$3.2\cdot 10^{7}$ & 0.144592 & 0.148365 & 0.149341 & 0.156702 & 0.161709 & 0.162451 & 0.182390 & 0.183691 \\
$6.4\cdot 10^{7}$ & 0.144593 & 0.148366 & 0.149341 & 0.156702 & 0.161709 & 0.162451 & 0.182390 & 0.183692 \\
$1.28\cdot 10^{8}$ & 0.144598 & 0.148368 & 0.149342 & 0.156703 & 0.161709 & 0.162451 & 0.182390 & 0.183691 \\
$2.56\cdot 10^{8}$ & 0.144598 & 0.148368 & 0.149343 & 0.156703 & 0.161710 & 0.162451 & 0.182390 & 0.183691 \\
\bottomrule
\end{tabular}

\caption{$f(k)=2k+1$.}
\end{table}
\clearpage

\section{Provenance}\label{app:provenance}

The mathematics in \cref{sec:prelim,sec:densities,sec:identity,sec:pairs,sec:sawtooth,sec:single,sec:poly,sec:assembly} is an edited presentation of an archived run of a multi-agent proof pipeline in which language models proposed, proved and refereed the results. This appendix records where the material comes from. A proposition-by-proposition map, with every editorial change and correction, is in \texttt{manuscript/changes.md}.

\emph{Archive.}
\begin{itemize}
\item Source repository: \url{https://gitlab.com/mshepherd/squarefree-numbers}, archive label \texttt{zt-neutral}.
\item Archive commit: \texttt{ed0aebaf461b602601856bf1b693fd76893d9aba}.
\item Harness commit: \texttt{33aadc609359221ba5d81664438bce5f658f561c}.
\item Revision recorded for all run stages: \texttt{8dfc53d01e56c9949c42061fdbe907a98a12a5d7}.
\end{itemize}

\emph{Models and roles (as recorded).} The model identifiers are those recorded in the archive. The command-line tool versions are not model-weight versions, and model versions beyond these identifiers are unknown.
\begin{itemize}[before=\raggedright]
\item Planning and probing: \texttt{claude-opus-5-5} and \texttt{codex:gpt-6.1-sol}; probing also \texttt{claude-sonnet-5-5}.
\item Numerics: \texttt{claude-sonnet-5-5} and \texttt{codex:gpt-6.1-sol} produced candidate computations. The \texttt{codex:gpt-6.1-sol} snapshot was accepted.
\item Proofs: \texttt{claude-sonnet-5-5} (prover).
\item Refereeing: \texttt{claude-opus-5-5} and \texttt{codex:gpt-6.1-sol}. Each proposition was reviewed independently by both.
\item Integration (assembly audit): \texttt{claude-opus-5-5} and \texttt{codex:gpt-6.1-sol}, independently.
\end{itemize}
Claude models ran under Claude Code 2.1.309 and Codex models under codex-cli 0.166.0.

\emph{Write-up stage.} This text was written by \texttt{claude-opus-5-5} (Claude Code 2.1.311) from the archived run. It was refereed by \texttt{claude-opus-5-5} and \texttt{codex:gpt-6.1-sol} (the write-up stage records Claude Code 2.1.311 and codex-cli 0.166.0) in three rounds. Each round was followed by a revision, and the present version answers the third round. Some of the corrections described below were raised by these write-up referees.

\emph{Workflow.} The problem statement restricts $f$ to $\Z[t]$. It asks for the asymptotic behaviour of $V_f(x,y)$ in as large a range of $f(y)\le x$ as can be justified, and it accepts a weaker range or error term provided it is correct. The weight $f'(k)$ and the class~\eqref{eq:f-hyp} of polynomials used here are those of the problem statement; neither was changed. The numerical candidates were computed before any plan existed, and one was accepted on computational grounds only, not because of its fitted constants. Eight proof plans were written, and the plan \texttt{simplest@claude-opus-5-5} was selected by the project owner. It splits the argument into nine propositions P1--P9. In this text, P1 is \cref{prop:densities}, P2 is \cref{prop:identity}, P3 is \cref{prop:pairs}, P4 is \cref{prop:trivial}, P5 is \cref{prop:sawtooth}, P6 is \cref{prop:MT}, P7 is \cref{prop:single}, P8 is \cref{prop:poly}, and P9 is \cref{thm:main} with its proof in \cref{sec:assembly}. All nine proofs passed both referees in the first round, with only cosmetic remarks. Both integrators found that the propositions compose to the stated theorem without loss and without interface mismatches. The project owner's recorded decision reads: ``Let's just assume the proof is correct and complete.'' A separate read-only operator audit found no substantive gap and left final mathematical acceptance to the project owner. The corrections made in this text are local and fall into three groups, listed individually in \texttt{changes.md}: (a) those recorded in the operator audit (the normalisation and convergence wording in P1(iii), the endpoint $k=1$ in P8, and unused numerical remarks in P6 and P8); (b) cosmetic points raised by the run's referees; and (c) points raised at the write-up stage. Group~(c) contains one clarification of an archived auxiliary statement. For an infinite interval, the Euler--Maclaurin lemma in the archived P5 proof says that ``the same holds'' under the hypotheses $G,G'\to0$, $G''\in L^1$. Read as an identity between a convergent series and a convergent integral, this is false in that generality. \Cref{lem:EM} keeps exactly the archived hypotheses and states what letting $b\to\infty$ gives: the difference between the partial sums and the partial integrals converges, and the formula holds whenever the series or the integral converges. Both archived applications have absolutely convergent series and integrals, so no proposition and no step of the argument changes. No decision of the project owner on this point is recorded, and \texttt{changes.md} lists it for the owner's attention. The other items in group~(c) concern only the exposition of this text, for example an arithmetic slip in an earlier version of the write-up, and do not touch the archived arguments. Agreement between models is evidence, not certification; the arguments are reproduced here in full so that they can be checked directly.

\emph{Independent checks during the run.} Before acceptance, the operator compared both numerical candidates on overlapping rows and recomputed some rows by brute force. One integrator recomputed all $16$ rows for $k^2+1$ and $k^3+2$ at $x=250000$ by direct counting. The other computed $\sum f'(k)\MT(x;f(k))$ independently and compared it with the data. These checks are recorded in the archive. The computations in \cref{sec:numerics} were redone for this write-up as described there.

\emph{Literature.} The run used the papers supplied with the problem and literature up to 2010, with the single exception of Brüdern--Wooley~\cite{BruedernWooley2011}, which was supplied. The present text cites only works in the supplied corpus. See \texttt{README.md} for the limitations of the bibliography.


\end{document}
